\documentclass[titlepage,a4paper,11pt]{article}

\usepackage[right=3cm, left=3cm, top=3cm, bottom=3cm]{geometry}

\usepackage{amsfonts}
\usepackage{amsthm}
\usepackage[latin1]{inputenc}
%\usepackage[T1]{fontenc}
\usepackage{euscript}
\usepackage{amssymb, amsmath, mathtools, enumerate, color, graphicx, mathrsfs}
%\usepackage{mathtools}
\usepackage{hyperref}%damit Referenzen springbar sind
\usepackage{stmaryrd} % f�r Widerspruchzeichen \lightning
\usepackage{paralist} % added by Ch�
%\usepackage{ulem} % f�r unterstreichen mit wellen
\usepackage{tikz}

% STYLE
\numberwithin{equation}{section}
\DeclareMathAlphabet{\mathpzc}{OT1}{pzc}{m}{it}
\hypersetup{colorlinks=true}
\bibliographystyle{alpha}
%\parindent 0pt% keine Abs�tze mehr


% THEOREMS
\newtheorem{theo}{Theorem}[section]
\newtheorem{kor}[theo]{Corollary}
\newtheorem{lemma}[theo]{Lemma}
\newtheorem{prop}[theo]{Proposition}

\theoremstyle{definition} % Definitionen nichtmehr kursiv geschrieben
\newtheorem{bsp}[theo]{Example}
\newtheorem{bem}[theo]{Remark}
\newtheorem{rem}[theo]{Remark}
\newtheorem{definition}[theo]{Definition}


%--------------------------------------------------------------------------------------------------
%
%     ANTON'S DEFINITIONS
%
%--------------------------------------------------------------------------------------------------
\renewcommand{\labelitemi}{} % keine Punkte bei itemize
\newcommand{\eq}{\Leftrightarrow}
\newcommand{\Ra}{\Rightarrow}
\newcommand{\La}{\Leftarrow}
\newcommand{\nin}{{n\in\N}}
\newcommand{\R}{\mathbb{R}}
\newcommand{\N}{\mathbb{N}}
\newcommand{\C}{\mathbb{C}}
\newcommand{\K}{\mathbb{K}}
\newcommand{\Q}{\mathbb{Q}}
\newcommand{\ph}{\varphi}
\newcommand{\eps}{\varepsilon}
\newcommand{\sub}{\subseteq}
\newcommand{\p}{\vspace{0.3cm}\\}
%--------------------------------------------------------------------------------------------------


%--------------------------------------------------------------------------------------------------
%
%     FRITZ' DEFINITIONS
%
%--------------------------------------------------------------------------------------------------

% Things I cannot remember
\newcommand{\paragraf}{\textsection}
\renewcommand{\emptyset}{\varnothing}



% Inner products
\newcommand{\lk}{\langle}
\newcommand{\rk}{\rangle}
\newcommand{\aproduct}{\langle\cdot\,,\cdot\rangle}
%\newcommand{\llb}{\llbracket}               % Nur mit package stmaryrd.sty
%\newcommand{\rrb}{\rrbracket}               % Nur mit package stmaryrd.sty
%\newcommand{\dproduct}{\llb\cdot,\cdot\rrb} % Nur mit package stmaryrd.sty
%\newcommand{\dperp}{{\llb\perp\rrb}}        % Nur mit package stmaryrd.sty


% Letters
\newcommand{\calA}{\mathcal A}         \newcommand{\frakA}{\mathfrak A}
\newcommand{\calB}{\mathcal B}         \newcommand{\frakB}{\mathfrak B}
\newcommand{\calC}{\mathcal C}         \newcommand{\frakC}{\mathfrak C}
\newcommand{\calD}{\mathcal D}         \newcommand{\frakD}{\mathfrak D}
\newcommand{\calE}{\mathcal E}         \newcommand{\frakE}{\mathfrak E}
\newcommand{\calF}{\mathcal F}         \newcommand{\frakF}{\mathfrak F}
\newcommand{\calG}{\mathcal G}         \newcommand{\frakG}{\mathfrak G}
\newcommand{\calH}{\mathcal H}         \newcommand{\frakH}{\mathfrak H}
\newcommand{\calI}{\mathcal I}         \newcommand{\frakI}{\mathfrak I}
\newcommand{\calJ}{\mathcal J}         \newcommand{\frakJ}{\mathfrak J}
\newcommand{\calK}{\mathcal K}         \newcommand{\frakK}{\mathfrak K}
\newcommand{\calL}{\mathcal L}         \newcommand{\frakL}{\mathfrak L}
\newcommand{\calM}{\mathcal M}         \newcommand{\frakM}{\mathfrak M}
\newcommand{\calN}{\mathcal N}         \newcommand{\frakN}{\mathfrak N}
\newcommand{\calO}{\mathcal O}         \newcommand{\frakO}{\mathfrak O}
\newcommand{\calP}{\mathcal P}         \newcommand{\frakP}{\mathfrak P}
\newcommand{\calQ}{\mathcal Q}         \newcommand{\frakQ}{\mathfrak Q}
\newcommand{\calR}{\mathcal R}         \newcommand{\frakR}{\mathfrak R}
\newcommand{\calS}{\mathcal S}         \newcommand{\frakS}{\mathfrak S}
\newcommand{\calT}{\mathcal T}         \newcommand{\frakT}{\mathfrak T}
\newcommand{\calU}{\mathcal U}         \newcommand{\frakU}{\mathfrak U}
\newcommand{\calV}{\mathcal V}         \newcommand{\frakV}{\mathfrak V}
\newcommand{\calW}{\mathcal W}         \newcommand{\frakW}{\mathfrak W}
\newcommand{\calX}{\mathcal X}         \newcommand{\frakX}{\mathfrak X}
\newcommand{\calY}{\mathcal Y}         \newcommand{\frakY}{\mathfrak Y}
\newcommand{\calZ}{\mathcal Z}         \newcommand{\frakZ}{\mathfrak Z}

\newcommand{\scrA}{\mathscr A}  % For scr-font package "mathrsfs" is needed
\newcommand{\scrB}{\mathscr B}
\newcommand{\scrC}{\mathscr C}
\newcommand{\scrD}{\mathscr D}
\newcommand{\scrE}{\mathscr E}
\newcommand{\scrF}{\mathscr F}
\newcommand{\scrG}{\mathscr G}
\newcommand{\scrH}{\mathscr H}
\newcommand{\scrI}{\mathscr I}
\newcommand{\scrJ}{\mathscr J}
\newcommand{\scrK}{\mathscr K}
\newcommand{\scrL}{\mathscr L}
\newcommand{\scrM}{\mathscr M}
\newcommand{\scrN}{\mathscr N}
\newcommand{\scrO}{\mathscr O}
\newcommand{\scrP}{\mathscr P}
\newcommand{\scrQ}{\mathscr Q}
\newcommand{\scrR}{\mathscr R}
\newcommand{\scrS}{\mathscr S}
\newcommand{\scrT}{\mathscr T}
\newcommand{\scrU}{\mathscr U}
\newcommand{\scrV}{\mathscr V}
\newcommand{\scrW}{\mathscr W}
\newcommand{\scrX}{\mathscr X}
\newcommand{\scrY}{\mathscr Y}
\newcommand{\scrZ}{\mathscr Z}

\newcommand{\la}{\lambda}
\newcommand{\veps}{\varepsilon}
\newcommand{\vphi}{\varphi}


% Matrices and vectors
\newcommand{\mat}[4]
{
   \begin{pmatrix}
      #1 & #2\\
      #3 & #4
   \end{pmatrix}
}
\newcommand{\vek}[2]
{
   \begin{pmatrix}
      #1\\
      #2
   \end{pmatrix}
}
\newcommand{\smallvek}[2]{\left(\begin{smallmatrix}#1\\#2\end{smallmatrix}\right)}
\newcommand{\smallmat}[4]{\left(\begin{smallmatrix}#1 & #2\\#3 & #4\end{smallmatrix}\right)}


% Not in italic mode
\renewcommand{\Im}{\operatorname{Im}}
\renewcommand{\Re}{\operatorname{Re}}
\DeclareMathOperator*{\linspan}{span}
\renewcommand{\ker}{\operatorname{ker}}
\DeclareMathOperator{\ran}{ran}
\DeclareMathOperator{\dom}{dom}
\DeclareMathOperator{\codim}{codim}


% Arrows
\newcommand{\lra}{\longrightarrow}
\newcommand{\sra}{\rightarrow}
\newcommand{\Lra}{\Longrightarrow}
\newcommand{\Sra}{\Rightarrow}

\newcommand{\lla}{\longleftarrow}
\newcommand{\sla}{\leftarrow}
\newcommand{\Lla}{\Longleftarrow}
\newcommand{\Sla}{\Leftarrow}

\newcommand{\llra}{\longleftrightarrow}
\newcommand{\slra}{\leftrightarrow}
\newcommand{\Llra}{\Longleftrightarrow}
\newcommand{\Slra}{\Leftrightarrow}

\newcommand{\upto}{\uparrow}
\newcommand{\downto}{\downarrow}

\newcommand{\wto}{\rightharpoonup}
\newcommand{\Ato}{\stackrel{A}{\to}}
\newcommand{\Awto}{\stackrel{A}{\rightharpoonup}}
%\newcommand{\restr}{\!\!\upharpoonright\!\!}


% Misc
\newcommand{\ol}{\overline}
\newcommand{\ds}{\dotplus}
\newcommand{\wt}{\widetilde}
\newcommand{\wh}{\widehat}

% Only here
\newcommand{\sse}{\subseteq}
%---------------------------------------------------------------------------------------------------------
%---------------------------------------------------------------------------------------------------------

% newcommand: Alex
\newcommand{\sequ}[2][n]{(#2_{#1})_{#1 \in \N}}				% short for sequences: Argument is the sequ element: optional index (default is n)
\newcommand{\serie}[2][1]{\sum_{#2=#1}^{\infty}}			% serie with given index and optional starting value

\newcommand{\inter}[1]{\stackrel{\circ}{#1}}				% the interior of a set
\newcommand{\closure}[2][3]{%								% closure of a set. 1 optional agument for litte costumization 
	{\mkern#1mu\overline{\mkern-#1mu#2}}}

\newcommand{\epsi}{\varepsilon > 0}							% epsi := e > 0, just a shortcut
\newcommand{\reci}[1]{\frac{1}{#1}}							% reciprocal
\newcommand{\Span}[1]{\operatorname{span\left\{#1\right\}}} % span over the argument
% END: newcommand: Alex


% newcommand: Che
\newcommand{\ub}{\underbrace}
\newcommand{\id}{\operatorname{\id}}

% END: newcommand: Che

%--------------------------------------------------------------------------------------------------
%
%     AXEL'S DEFINITIONS
%
%--------------------------------------------------------------------------------------------------

\newcommand{\norm}[1]{\Vert #1 \Vert}
\newcommand{\dual}[1]{{#1}^{\ast}}
\newcommand{\bidual}[1]{{#1}^{\ast \ast}}

\newcommand{\sequence}[2]{\ensuremath{(#1)_{#2 \in \mathbb{N}}}}

% Martins Commands
\newcommand{\field}[1]{\mathbb{#1}}


%\newcommand{\N}{\field{N}}
%\newcommand{\R}{\field{R}}
%\newcommand{\C}{\field{C}}
%\newcommand{\H}{\field{H}}
\newcommand{\Z}{\field{Z}}
%\newcommand{\Q}{\field{Q}}
%\newcommand{\K}{\field{K}}

\newcommand{\KK}{\mathcal{K}}
\newcommand{\RR}{\mathcal R}
\newcommand{\M}{\mathcal M}
\newcommand{\F}{\mathcal F}
\newcommand{\B}{\mathcal B}

\renewcommand{\P}{\mathcal P}
\renewcommand{\i}{\mathrm{i}\,}
\renewcommand{\L}{\mathscr{L}}


%\newcommand{\eps}{\varepsilon}
%\newcommand{\ph}{\varphi}

\newcommand{\const}{\mathrm{const}}
%\newcommand{\ran}{\mathrm{Ran}}
\newcommand{\phys}{\mathrm{ph}}
\newcommand{\modi}{\mathrm{mod}}
\newcommand{\Res}{\mathrm{Res}}
\newcommand{\SO}{\mathrm{SO}}
\newcommand{\interior}[1]{\mbox{$\overset{\circ}{#1}$}}
\newcommand{\e}{\mathrm{e}}
\newcommand{\abl}{\medspace\mathrm{d}}
\newcommand{\lap}{\enspace\laplace\enspace}
\newcommand{\Lap}{\enspace\Laplace\enspace}
\newcommand{\re}{\mathrm{Re}}
\newcommand{\im}{\mathrm{Im}}
\newcommand{\disp}{\displaystyle}

\newcommand{\Ran}{\operatorname{Ran}}
\newcommand{\supp}{\operatorname{supp}}
\newcommand{\Ima}{\operatorname{Im}}
\newcommand{\Rea}{\operatorname{Re}}
\newcommand{\dime}{\operatorname{dim}}
\newcommand{\card}{\operatorname{card}}
\renewcommand{\div}{\operatorname{div}}
\newcommand{\rot}{\operatorname{rot}}
\newcommand{\spur}{\operatorname{Sp}}
\newcommand{\dist}{\operatorname{dist}}

\newcommand{\restricted}[1]{|\grave{#1}\,}
\newcommand{\abs}[1]{\left|#1\right|}

\newcommand{\set}[1]{\left\lbrace #1 \right\rbrace}
%\renewcommand{\span}[1]{\left\langle #1 \right\rangle}
\newcommand{\thf}{\Rightarrow}
\renewcommand{\iff}{\Leftrightarrow}

\newcommand{\contradiction}{%
\begin{tikzpicture}[rotate=45,x=0.5ex,y=0.5ex]
\draw[line width=.2ex] (0,2) -- (3,2) (0,1) -- (3,1) (1,3) -- (1,0) (2,3) -- (2,0);
\end{tikzpicture}
}


%\newcommand{\sequence}[2]{\ensuremath{(#1)_{#2 \in \mathbb{N}}}}


\begin{document}

\begin{titlepage}
\title{{\bf Functional Analysis 1}\\[2cm]\uppercase{\small {Lecture Script}}}
\author{Prof. Dr. Gitta Kutyniok\\[10cm]}
\date{\today}
\maketitle
\end{titlepage}
\section{Metric spaces}

In this chapter we recall the basic notions of metric spaces and prove Baire's theorem and the theorem of Arzela-Ascoli. Throughout this lecture $\K$ will always denote either $\R$ or $\C$.

\begin{definition}
	Let $X$ be a set. Then a map $d\colon X\times X\to[0,\infty)$ is called a \emph{metric on $X$}, if for all $x,y,z\in X$	
	\begin{align*}
	\hspace{-1cm}&(i)  \hspace{-2cm}&d(x,y)&=0 \eq x=y.\\
	\hspace{-1cm}&(ii)  \hspace{-2cm}&d(x,y)&=d(y,x).\\
	\hspace{-1cm}&(iii)  \hspace{-2cm}&d(x,z)&\leq d(x,y)+d(y,z). \tag{Triangle inequality}
	\end{align*}
	$(X,d)$ is then called a \emph{metric space}, and $d(x,y)$ is referred to as the \emph{distance} between $x$ and $y$.
	If $Y\subseteq X$, then $d|_{Y\times Y}$ is the \emph{induced metric on $Y$}.
\end{definition}

Notice that the non-negativity of a metric already follows from
\[
	0 = d(x,x) \leq d(x,y)+d(y,x) = 2d(x,y).
\]
Next we will give some important examples of metrics on function spaces, sequence spaces and $\K^n$. Also, we can define a metric on every set as the first example will show.

\begin{bsp}\label{lp-spaces}~\begin{enumerate}
	\item[(1)]
		Let $X$ be a set and let $d\colon X\times X\to [0,\infty)$ be defined by
		\begin{align*}
			d(x,y):=\begin{cases}1:x\neq y\\0:\text{else}\end{cases}
		\end{align*}
		This is the so-called \emph{discrete metric}. Hence this always defines a metric.
	\item[(2)]
		Let $X$ be a set and define
		\begin{align*}
			B(X)=\{f\colon X\to \K:f\text{ is bounded}\}.
		\end{align*}
		Then
		\begin{align*}
			d(f,g):=\sup\limits_{x\in X} |f(x)-g(x)|
		\end{align*}
		is a metric on $B(X)$, the so-called \emph{supremum metric}.
		Let now $X=[a,b]$ and set
		\begin{align*}
			C[a,b]=\{f\colon[a,b]\to\K:f\text{ continuous}\}.
		\end{align*}
		Then
		\begin{align*}
			C[a,b]\subseteq B[a,b]
		\end{align*}
		and hence $d$ induces a metric on $C[a,b]$.
	\item[(3)]
		For $1\leq p<\infty$, let $d_p\colon\K^n\times\K^n\to[0,\infty)$ be defined by
		\begin{align*}
			d_p(x,y)&:= \left(\sum\limits_{j=1}^n|x_j-y_j|^p\right)^\frac1p,~x=(x_j)_{j\in\N},~y=(y_j)_{j\in\N},
		\intertext{and let $d_\infty\colon\K^n\times\K^n\to[0,\infty)$ be defined by}
			d_\infty(x,y)&:=\max\limits_{1\leq j\leq n}|x_j-y_j|.
		\end{align*}
		$p=1,\infty$ : This is obviously a metric\\
		$1<p<\infty$ : Theorem \ref{minkowski} will imply that this is a metric.
	\item[(4)]
		The spaces $(\K^n,d_p)$ can be generalized to "infinite-dimensional sequence spaces". For this, for $1\leq p<\infty$, set
		\begin{align*}
			\ell_p:=\left\{x=(x_n)_\nin:x_n\in\K,~\sum\limits_{n=1}^\infty|x_n|^p<\infty\right\},
		\end{align*}
		and define $d_p\colon \ell_p\times \ell_p\to[0,\infty)$ by
		\begin{align*}
			d_p(x,y)=\left(\sum\limits_{n=1}^\infty|x_n-y_n|^p\right)^\frac1p.
		\end{align*}
		This is well-defined, since by Theorem \ref{minkowski} $\ell_p$ is a linear space. Let further $\ell_\infty$ be defined by
		\begin{align*}
			\ell_\infty:=\left\{ x=(x_n)_\nin:(x_n)_\nin\text{ bounded}\right\}
		\intertext{and define $d_\infty\colon \ell_\infty\times \ell_\infty\to[0,\infty)$ by}
			d_\infty(x,y):=\sup\limits_\nin|x_n-y_n|.
		\end{align*}
		Then $(\ell_p,d_p)$, $1\leq p\leq\infty$ are metric spaces, again partly proven by Theorem \ref{minkowski}.
	\end{enumerate}
\end{bsp}

To show the triangle inequality for the $\ell_p$-spaces we need another inequality, which is important in its own right. H�lder's inequality gives upper bounds on a series of products in terms of products of series.

\begin{theo}[H�lder's inequality]\label{holder}
	Let $1<p<\infty$, and let $1<q<\infty$ be defined by $q:=\frac{p}{p-1}$ (hence $\frac1p+\frac1q=1$). 
	Then, for $x\in \ell_p$ and $y\in \ell_q$, we have
	\begin{align*}
		\sum\limits_{n=1}^\infty|x_n {y}_n|
		\leq\left(\sum\limits_{n=1}^\infty|x_n|^p\right)^\frac1p\left(\sum\limits_{n=1}^\infty|y_n|^q\right)^\frac1q.
	\end{align*}
	(Case $p=q=2$: Schwarz' inequality)
\end{theo}
\begin{proof}
	Let $c=\frac1p$ and define $\ph\colon[0,\infty)\to\R$ by $\ph(t)=t^c-ct$. Then
	\begin{align*}
		\ph'(t)=ct^{c-1}-c~\text{and}~\ph''(t)=c(c-1)t^{c-2}.
	\end{align*}
	Thus $\ph$ has a global maximum value in $t=1$. This implies
	\begin{align*}
		1-c&\geq t^c-ct~\text{for all}~t>0,
	\intertext{hence}
		t^c-1&\leq c(t-1)\tag{$\ast$}.\label{ast}
	\intertext{Let now $a,b>0$, and set $t=\frac{a^p}{b^q}$. Then, by \eqref{ast}, we obtain}
		\frac{a}{b^{\frac{q}{p}}}-1\leq \frac1p\left(\frac{a^p}{b^q}-1\right)
		&\Ra \frac{a}{b^{q\left(\frac1p-1\right)}}-b^q\leq\frac1p\left(a^p-b^q\right).
	\intertext{Since $1=\frac1p+\frac1q$, this implies}
		ab&\leq\frac{a^p}{p} +\frac{b^q}{q} .\tag{$\ast\ast$}\label{astast}
	\end{align*}
	We now set
	\begin{align*}
		A:=\left(\sum\limits_{n=1}^\infty|x_n|^p\right)^\frac1p~\text{and}~
		B:=\left(\sum\limits_{n=1}^\infty|y_n|^q\right)^\frac1q
	\end{align*}
	as well as $\widetilde x_n:=\frac{x_n}A$ and $\widetilde y_n:=\frac{y_n}B$. WLOG, we assume $A,B>0$. By \eqref{astast}, we obtain
	\begin{align*}
		|\widetilde x_n\widetilde y_n|&\leq \frac1p|\widetilde x_n|^p+\frac1q|\widetilde y_n|^q.
	\intertext{Hence}
		\sum\limits_{n=1}^\infty|\widetilde x_n \widetilde y_n|
		&\leq\frac1p\sum\limits_{n=1}^\infty|\widetilde x_n|^p+\frac1q\sum\limits_{n=1}^\infty|\widetilde y_n|^q
		=\frac1p+\frac1q=1.
	\intertext{And, finally:}
		\sum\limits_{n=1}^\infty|x_n y_n|&\leq AB,
	\end{align*}
which is the assertion.
\end{proof}

The following Minkowski's inequality sets the ground for the triangle inequality of the metric $d_p$.

\begin{theo}[Minkowski's inequality]\label{minkowski}
	For $1<p<\infty$ and $x,y\in \ell_p$,
	\begin{align*}
		\left(\sum\limits_{n=1}^\infty |x_n+y_n|^p\right)^\frac1p
		\leq\left(\sum\limits_{n=1}^\infty |x_n|^p\right)^\frac1p+\left(\sum\limits_{n=1}^\infty |y_n|^p\right)^\frac1p.
	\end{align*}
\end{theo}
\begin{proof}
	With $z_n:=x_n+y_n$, we first obtain
	\begin{align*}
		|z_n|^p&\leq\sum\limits_{n=1}^m|x_n+y_n|\cdot|z_n|^{p-1}\leq\left(|x_n|+|y_n|\right)|z_n|^{p-1}.
	\intertext{This implies}
		\sum\limits_{n=1}^m|z_n|^p&\leq \sum\limits_{n=1}^m|x_n|\cdot|z_n|^{p-1}+\sum\limits_{n=1}^m|y_n|\cdot|z_n|^{p-1}
		~\text{for all}~m\in\N.
	\intertext{By Theorem \ref{holder},}
		\sum\limits_{n=1}^m|z_n|^p&\leq\left(\sum\limits_{n=1}^m|x_n|^p\right)^\frac1p \left(\sum\limits_{n=1}^m|z_n|^{(p-1)q}\right)^\frac1q
			+\left(\sum\limits_{n=1}^m|y_n|^p\right)^\frac1p \left(\sum\limits_{n=1}^m|z_n|^{(p-1)q}\right)^\frac1q.
	\end{align*}
	Since $(p-1)q=p$, we conclude that
	\begin{align*}
		\left(\sum\limits_{n=1}^m|z_n|^p\right)^\frac1p\left(\sum\limits_{n=1}^m|z_n|^p\right)^{1-\frac1q}
			\leq\left(\sum\limits_{n=1}^m|x_n|^p\right)^\frac1p+\left(\sum\limits_{n=1}^m|y_n|^p\right)^\frac1p.
	\end{align*}
	We now consider $m\to\infty$, which we are allowed to do since the right-hand-side converges. This proves the theorem.
\end{proof}

The triangle inequality $d_p(u,w)\leq d_p(u,v)+d_p(v,w)$ for all $u,v,w\in \ell_p$ for the metric $d_p$ can now directly be concluded from Theorem \ref{minkowski} by setting $x_n=u_n-v_n$ and $y_n=v_n-w_n$.

\begin{definition} Let $(X,d)$ be a metric space.
\begin{enumerate}
	\item[(1)] For $x\in X$ and $r>0$, the set $U_r(x)$ defined by
		\begin{align*}
			U_r(x):=\{y\in X:d(x,y)<r\}
		\end{align*}
		is called the \emph{open ball} of radius $r$ and center $x$. $U\sub X$ is called \emph{open}, if for each
		$x\in U$ there exists some $\eps>0$ such that $U_\eps(x)\sub U$.
	\item[(2)] A set $A\sub X$ is \emph{closed}, if $X\backslash A$ is open. The set
		\begin{align*}
			K_r(x):=\{y\in X:d(x,y)\leq r\},~x\in X,~r>0,
		\end{align*}
		is called the \emph{closed ball} of radius $r$ and center $x$.
	\item[(3)] If $E\subseteq X$, then $x\in E$ is an \emph{interior point of $E$}, if there exists some open set $U\subseteq X$
		with $x\in U\subseteq E$. $E$ is then called a \emph{neighbourhood of $x$}. The set of all interior points is refered to as the
		\emph{interior of $E$} and is denoted by $\mathring E$.
	\item[(4)] A point $x\in X$ is called \emph{limit point of E} if $U\cap E\neq\emptyset$ for each neighbourhood $U$ of $x$.
		The set of all limit points of $E$ is the \emph{closure of $E$}, which is denoted by $\overline E$. $E$ is \emph{dense} in $X$ if 
		$\overline E=X$.
\end{enumerate}
\end{definition}

The openness of a set and all properties that can be defined with reference only to open sets are called topological. In particular, all terms just defined are topological. The open sets in a metric space form a system of sets called topology. 
This terminology will be generalized to the notion of a topological space in Chapter X\marginpar{Reference!}.

\begin{lemma} Let $(X,d)$ be a metric space.
\begin{enumerate}
	\item[(i)] We have \begin{enumerate}
		\item[(a)] $\emptyset$, $X$ are open.
		\item[(b)] $U_1,\ldots U_r\subseteq X$ open $\Ra$ $\bigcap\limits_{i=1}^rU_i$ is open.
		\item[(c)] $U_i\subseteq X$, $i\in I$ open $\Ra$ $\bigcup\limits_{i\in I}U_i$ is open.
		\end{enumerate}
		Hence $d$ defines a \emph{topology} on $X$ with $U_\eps(x)$, $x\in X$, $\eps>0$, as basis.
	\item[(ii)] We have \begin{enumerate}
		\item[(a)] $\emptyset$, $X$ are closed.
		\item[(b)] $A_i\subseteq X$, $i\in I$ closed $\Ra$ $\bigcap\limits_{i\in I}A_i$ is closed.
		\item[(c)] $A_1,\ldots A_r\subseteq X$ open $\Ra$ $\bigcup\limits_{i=1}^rA_i$ is closed.
		\end{enumerate}
	\item[(iii)] For each $x\in X$, $r>0$, the set $K_r(x)$ is closed.
	\item[(iv)] For $E\subseteq X$, $\overline E$ is the smallest closed set containing $E$.
	\item[(v)] For $E\subseteq X$, $\mathring E$ is the biggest open set contained in $E$.
\end{enumerate}
\end{lemma}
\begin{proof} Tutorials\end{proof}

The next definition generalizes the notion of convergence from $\K^n$ (with the Euclidian metric) to general metric spaces.

\begin{definition} Let $(X,d)$ be a metric space.
\begin{enumerate}
	\item[(1)] A sequence $(x_n)_\nin\subseteq X$ \emph{converges} to $x\in X$ if for each $\eps>0$ there exists
		$N_\eps\in\N$ with $d(x_n,x)<\eps$ for all $n\geq N_\eps$. We then write $x_n\to x$, as $n\to\infty$ or 
		$x=\lim\limits_{n\to\infty}x_n$. $x$ is called the \emph{limit} of $(x_n)_\nin$.
	\item[(2)] A sequence $(x_n)_\nin\subseteq X$ is a \emph{Cauchy-sequence}, if for each $\eps>0$ there exists some
		$N_\eps\in\N$ with
		\begin{align*} d(x_n,x_m)<\eps\quad\text{for all }n,m\geq N_\eps .\end{align*}
	\item[(3)] $(X,d)$ is \emph{complete}, if each Cauchy-sequence in $X$ converges.
\end{enumerate}
\end{definition}

Convergence of a sequence is a topological property. The sequence $(x_n)_\nin$ converges to $x$, if and only if every neighborhood of $x$ contains all but finitely many elements of the sequence. In particular, the limit of a sequence is independent of the ordering of the sequence's terms. Which sequences are Cauchy-sequences does not only depend on the open sets but also on the chosen metric (see Remark \ref{completeness is not topological}).

In general, topological properties in metric spaces can be tested by sequences. We note that there is a characterization of closedness of a set by convergent sequences.

\begin{lemma}\label{18} Let $(X,d)$ be a metric space.
\begin{enumerate}
	\item[(i)] A sequence can have at most one limit.
	\item[(ii)] Let $E\sub X$. Then $x\in\overline E$ if and only if there exists a sequence $(x_n)_\nin\sub E$ with $x_n\to x$
		as $n\to\infty$.
	\item[(iii)] If $(x_n)_\nin\sub X$ is convergent, then $(x_n)_\nin\sub X$ is a Cauchy-sequence. 
		The converse is not always true\footnote{For example, consider $X=(0,1]$, $x_n=\frac 1n$.}. A Cauchy-sequence is convergent, if it contains a convergent sub\-sequence.
	\item[(iv)] If $X$ is complete and $E\sub X$ closed, then $E$ is complete. If $E\sub X$ is complete, then $E$ is closed in $X$.
\end{enumerate}
\end{lemma}
\begin{proof} Tutorials\end{proof}

The following example provides the reader with some complete metric spaces.

\begin{bsp}~
\begin{enumerate}
	\item[(1)] $B(X)$ is complete.
		\begin{proof}
			Let $(f_n)_\nin$ be a Cauchy-sequence in $B(X)$, and for $\eps>0$ let $N_\eps\in\N$ be such that
			\begin{align*}d(f_n,f_m)<\eps \text{ for all }n,m\geq N_\eps.\end{align*}
			This implies $|f_n(x)-f_m(x)|<\eps$ for all $x\in X$, $n,m\geq N_\eps$. Hence, for all $x\in X$, $(f_n(x))_\nin$
			is a Cauchy-sequence in $\K$. Setting $f(x):=\lim\limits_{n\to\infty}f_n(x)$, we obtain
			\begin{align*}|f(x)-f_m(x)|&=\lim\limits_{n\to\infty}|f_n(x)-f_m(x)|<\eps\quad\forall m\geq N_\eps.
				\intertext{Hence $|f(x)|\leq|f_n(x)|+\eps$, which implies $f\in B(X)$. Further, for $m\geq N_\eps$,}
				d(f,f_m)&=\sup\limits_{x\in X}|f(x)-f_m(x)|<\eps,
			\end{align*}
			and thus $f=\lim\limits_{n\to\infty}f_n$.
		\end{proof}
	\item[(2)] $C[a,b]$ is complete, since it is closed in $B[a,b]$ (see lemma \ref{18}), the reason being that a uniform 
		limit of continuous functions is again continuous.
	\item[(3)] $(\K^n,d_p)$, $n\in\N$, $1\leq p\leq\infty$ is complete, since convergence in $\K^n$ w.r.t. $d_p$ is the same as convergence
		in $\K^n$ w.r.t. the component sequences.
	\item[(4)] The spaces $\ell_p$, $1\leq p\leq \infty$ are complete.
		\begin{proof} Let $(x_k)_{k\in\N}$ be a Cauchy-sequence in $\ell_p$, $x_k=(x_{k,n})_\nin$, and for $\eps>0$ let $N_\eps\in\N$
		be with
		\begin{align*}
			d_p(x_k,x_l)&=\left(\sum\limits_{n=1}^\infty|x_{k,n}-x_{l,m}|^p\right)^\frac 1p<\eps&&\hspace{-2cm}
				\text{and}\label{stern}\tag{$\ast$}\\
			d_\infty(x_k,x_l)&=\sup\limits_\nin|x_{k,n}-x_{l,n}|<\eps&&\hspace{-2cm}\text{for all }k,l>N_\eps\label{sstern}\tag{$\ast\ast$}
		\end{align*}
		Thus, for fixed $n\in N_\eps$, $(x_{k,n})_{k\in\N}$ is a Cauchy-sequence in $\K$. Now set $y_n:=\lim\limits_{k\to\infty}x_{k,n}$
		and $y:=(y_n)_\nin$. Then $y\in \ell_p$ and $y=\lim\limits_{k\to\infty}x_n$.\\
		\underline{Reason:} Consider $l\to\infty$, which implies (by \eqref{stern},\eqref{sstern})
		\begin{align*}
			\sum\limits_{n=1}^m|x_{k,n}-y_n|^p&<\eps^p~\forall m\in\N
			\Ra \sum\limits_{n=1}^\infty |x_{k,n}-y_n|^p<\eps^p~\forall k\geq N_\eps\\
			\text{and }|x_{k,n}-y_n|&<\eps~\forall k\geq N_\eps,~n\in\N.
		\end{align*}
		Hence $x_n-y\in \ell_p$, and thus $y\in \ell_p$ and $y=\lim\limits_{k\to\infty}x_k$.
		\end{proof}
\end{enumerate}
\end{bsp}

The following theorem of Baire only holds in complete metric spaces. It is a key ingredient in the proofs of the fundamental theorems of functional analysis. Thus, they will only hold under some completeness assumption.

\begin{theo}[Baire's theorem]\label{baire}
	Let $(X,d)$ be a complete metric space, and let $D_n$, $\nin$ be open, dense subsets of $X$. Then also 
	$\bigcap\limits_\nin D_n$ is dense in $X$.
\end{theo}
\begin{proof}
	We need to prove that for all $x\in X$ and $r>0$ we have
	\begin{align*} U_r(x)\cap\bigcap\limits_{n=1}^\infty D_n\neq\emptyset.\end{align*}
	For this, let $x\in X$ and $r>0$ be arbitrary, but fixed. By induction define a sequence $(x_n)_\nin\sub X$ and
	$(r_n)_\nin\sub\R^+$ by
	\begin{enumerate}[a)]
		\item\label{blabli} $K_{r_{n+1}}(x_{n+1})\sub D_n\cap U_{r_n}(x_n)$
		\item\label{blibla} $r_n\leq\frac1n$
	\end{enumerate}
	This can be done as follows: First, set $x_1=x$ and $r_1=\min\{1,r\}$.
	Second, assume that $x_1,\ldots,x_n,r_1,\ldots,r_n$ be already chosen ($n\geq1$). Since $D_n$ is open and dense, also
	$D_n\cap U_{r_n}(x_n)\neq\emptyset$ is open. Hence there exists $x_{n+1}\in X$ and $r_{n+1}>0$ with
	\begin{align*}U_{2r_{n+1}}(x_{n+1})\sub D_n\cap U_{r_n}(x)\text{ and } r_{n+1}\leq\frac1{n+1}.\end{align*}
	This implies \ref{blabli}) and \ref{blibla}), since $K_{r_{n+1}}(x_{n+1})\sub U_{2r_{n+1}}(x_{n+1})$.
	\p
	Having constructed sequences $(x_n)$ and $(r_n)$ satisfying a) and b), we obtain
	\begin{align*}x_n\in K_{r_n}\sub D_{n-1}\cap U_{r_{n-1}}(x_{n-1})\sub U_{r_{n-1}}\sub\ldots\sub U_{r_m}(x_m)\end{align*}
	for all $n>m$. Thus $d(x_n,x_m)<r_m\leq\frac1m$ for all $n>m$. This implies that $(x_n)_\nin$ is a Cauchy-sequence in X.
	\p
	Now set $x_0:=\lim\limits_{n\to\infty}x_n$ (remember that $X$ is complete). Since $d(x_n,x_m)\leq r_m$ for all $n>m$, we obtain
	$d(x_0,x_m)\leq r_m$ for all $m\in\N$. Thus, finally,
	\begin{align*}
	x_0\in\bigcap\limits_{m=1}^\infty K_{r_{m+1}}(x_{m+1})\sub\bigcap\limits_{m=1}^\infty D_m\cap U_{r_m}(x_m)
	\sub U_{r_1}(x_1)\cap\bigcap\limits_{m=1}^\infty D_m\sub U_r(x)\cap\bigcap\limits_{m=1}^\infty D_m,
	\end{align*}
	and the theorem is proved.
\end{proof}

\begin{bem}~
	\label{rem:baire}
\begin{enumerate}
\item[(a)]	Theorem \ref{baire} is in general false if $X$ is not complete. As an example choose $X=\Q=\{q_1,q_2,\ldots\}$
	and $D_n=X\backslash\{q_n\}$, $n\in\N$, which are open and dense. We immediately see that however
	\begin{align*}\bigcap\limits_{n=1}^\infty D_n=\emptyset.\end{align*}
\item[(b)]	Let $(X,d)$ be complete and $A_n\sub X$, $n\in\N$ closed with $X=\bigcup\limits_{n=1}^\infty A_n$. 
	Then there exists at least one $n\in\N$ with
	\begin{align*} \mathring A_n\neq\emptyset.\end{align*}
	\begin{proof} Towards a contradiction, assume that
		\begin{align*}\mathring A_n=\emptyset\text{ for all }n\in\N.\end{align*}
		Then $X\backslash A_n$ are open and dense for all $n\in\N$. By Baire's theorem \ref{baire}, 
		$\bigcap\limits_{n=1}^\infty(X\backslash A_n)$ dense in $X$. But 
		$\bigcap\limits_{n=1}^\infty(X\backslash A_n)=X\backslash\bigcup\limits_{n=1}^\infty A_n=\emptyset$. $\lightning$
	\end{proof}
\item[(c)]\label{completeness is not topological} Completeness is a property of the particular metric and \underline{not} the convergence in $X$. For example, consider
	$X=(0,1]$, $d_1(x,y):=\left|\frac 1x-\frac1y\right|$ and $d_2(x,y)=|x-y|$. Then we have
	\begin{align*}x_n\to x\text{ in }(X,d_1)\eq x_n\to x\text{ in }(X,d_2),\end{align*}
	but $(X,d_1)$ is complete and $(X,d_2)$ is not (see tutorials).
\end{enumerate}
\end{bem}

\begin{definition} Let $(X,d)$ be a metric space.
\begin{enumerate}
	\item[(1)] Let $\eps>0$. Then $M\sub X$ is called \emph{$\eps$-net}, if $X=\bigcup\limits_{x\in M}U_\eps(x)$.
		$X$ is called \emph{totally bounded}, if for each $\eps>0$ there exists a finite $\eps$-net.
		$A\sub X$ is \emph{totally bounded}, if $(A,d|_{A\times A})$ is totally bounded.
	\item[(2)] $X$ is \emph{compact}, if every open cover of $X$ (that is, a family of open sets $U_i$, $i\in I$, such that $X=\bigcup_{i\in I}U_i$) has a finite subcover.
		($A,d|_{A\times A}$) is \emph{compact} if and only if every open cover of $A$ (of open sets in $X$) has a finite subcover.
\end{enumerate}
\end{definition}

Compactness and total boundedness are intrinsic properties, that is a subset $A \subseteq (X,d)$ of some metric space is compact (totally bounded) if the metric space $(A,d|_{A\times A})$ is compact (totally bounded).

It is easy to see that every compact metric space is totally bounded. The following theorem shows that the two notions coincide for complete metric spaces. Note that this does not imply that these two properties coincide for all subsets of a complete metric space (see Corollary \ref{cor_comp_bounded}).

\begin{theo}\label{bounded_compact} Let $(X,d)$ be a metric space. Then the following are equivalent:
\begin{enumerate}[(i)]
	\item\label{bounded}$(X,d)$ is complete and totally bounded.
	\item\label{compact}$(X,d)$ is compact.
	\item\label{conv_subseq} Each finite sequence in $X$ has a convergent subsequence.
\end{enumerate}
\end{theo}
\begin{proof}~
(\ref{bounded})$\Ra$(\ref{compact}). Towards a contradiction, assume that $X$ is not compact. 
	Let $\mathfrak{A}$ be an open cover of $X$ which does not
		contain a finite subcover. By induction, we now define a sequence $(x_n)_\nin\sub X$ satisfying
		\begin{enumerate}[a)]
			\item\label{finite}$U_{2^{-n}}(x_n)$ is not covered by finitely many $U\sub\mathfrak{A}$.
			\item\label{notdisjoint}$U_{2^{-n}}(x_n)\cap U_{2^{-(n+1)}}(x_{n-1})\neq\emptyset$.
		\end{enumerate}
		First, for $n=1$, notice that $X$ is totally bounded. Hence $X=\bigcup\limits_{y\in M}U_{\frac12}(y)$, $|M|<\infty$, 
		which implies that there exists $y_{i_0}=:x_1\in X$ such that $U_\frac12(x_1)$ is not covered by finitely many $U\sub\mathfrak{A}$.
		Second ($n\to n+1$), again by totally boundedness, there exists a finite $M$ such that 
		$X=\bigcap\limits_{y\in M}U_{2^{-(n+1)}}(y)$. Assume $x_1,\ldots,x_n$ are chosen such that \ref{finite}) and \ref{notdisjoint}) 
		are satisfied.
		Towards a contradiction assume that for each $y\in M$ with $U_{2^{-(n+1)}}(y)\cap U_{2^{-n}}(x_n)\neq\emptyset$, the set
		$U_{2^{-(n+1)}}(y)$ is covered by finitely many $U\in\mathfrak{A}$. Then this is also true for $U_{2^{-n}}(x_n)$ $\lightning$.
		Hence there exists $x_{n+1}\in X$ with $U_{2^{-(n+1)}}(x_{n+1})$ is not covered by finitely many $U\in\mathfrak{A}$ and
		$U_{2^-n}(x_n)\cap U_{2^{-(n+1)}}(x_{n+1})\neq\emptyset.$
		\p
		For each $n\in\N$, let $z_n\in U_{2^{-n}}(x_n)\cap U_{2^{-(n+1)}}(x_{n+1})$. Then, for $m>n$,
		\begin{align*}
			d(x_m,x_n)\leq\sum\limits_{\nu=n}^{m-1}d(x_{\nu+1},x_\nu)
			\leq&\sum\limits_{\nu=n}^{m-1}\left( d(x_{\nu+1},z_\nu)+d(z_\nu,x_\nu)\right)\\
			\leq&\sum\limits_{\nu=n}^{m-1}\left(2^{-(\nu+1)}+2^{-\nu}\right)\leq2\sum\limits_{\nu=n}^{m-1}2^{-\nu}\leq\frac{1}{2^{n-2}}.		
		\end{align*}
		This implies that $(x_n)_\nin$ is a Cauchy-sequence. Since $X$ is complete, there exists $x=\lim\limits_{n\to\infty}x_n$.
		\p
		Now choose $U\sub\mathfrak{A}$ with $x\in U$ and choose $\eps>0$ such that $U_\eps(x)\sub U$. Then
		$(x_n)\in U_{\frac{\eps}{2}}(x)$ for all $n\geq N$, hence $U_{2^{-n}}(x_n)\sub U$ for all $n\geq N$ with $2^{-n}<\frac\eps2$
		$\lightning$ to choice of $U_{2^{-n}}(x_n)$.
		
(\ref{compact})$\Ra$(\ref{conv_subseq}).
		Let $(x_n)_\nin$ be a sequence in $X$ and set $A_n:=\overline{\{x_\nu:\nu>n\}}\sub X$. Towards a contradiction assume that
		\begin{align*}\bigcap\limits_\nin A_n=\emptyset.\end{align*}
		This implies $\bigcup\limits_\nin(X\backslash A_n)=X.$ Since $X$ is compact, the open cover $\{X\backslash A_n\}_\nin$
		contains an open subcover $\{X\backslash A_{n_j}:1\leq j\leq r\}$. Since $A_{n+1}\sub A_n$, hence 
		$X\backslash A_n\sub X\backslash A_{n+1}$, for $N:=\max\{n_j:1\leq j\leq r\}$ we have
		\begin{align*} X=\bigcup\limits_{j=1}^rX\backslash A_{n_j}=X\backslash A_N.\end{align*}
		Thus $A_N=\emptyset$ $\lightning$.\\
		This proves $\bigcap\limits_\nin A_n\neq\emptyset$. Choosing $x\in\bigcap\limits_\nin A_n$, there exists a
		sequence $(n_k)_{k\in\N}\sub\N$ with $n_{k+1}>n_k$ and $d(x_{n_k},x)\leq\frac1k$ [if $n_k$ is chosen, then $x\in A_{n_{k+1}}$].
		This shows (\ref{conv_subseq}), since $(x_{n_k})_{k\in\N}$ is a convergent subsequence of $X$.

(\ref{conv_subseq})$\Ra$(\ref{bounded}). Each Cauchy-sequence in $X$ contains by hypothesis a convergent susequence, is hence itself convergent.
		This implies that $X$ is complete.\\
		Towards a contradiction, we now assume that $X$ is not totally bounded. Then there exists $\eps>0$ such that $X$ is not covered
		by finitely many $U_\eps(x)$, $x\in X$. By induction, we define a sequence $(x_n)_\nin\sub X$ with
		\begin{align*} x_n\notin U_\eps(x_j),~1\leq j\leq n-1.\end{align*}
		This can be achieved in the following way: Let $x_1\in X$ be arbitrary. Then assume $x_1,\ldots,x_n$ are already constructed.
		Since
		\begin{align*}
		X\backslash\bigcup\limits_{j=1}^n U_\eps(x_j)&\neq\emptyset,
		\intertext{choose $x_{n+1}\in X\backslash\bigcup\limits_{j=1}^n U_\eps(x_j).$ Then, for $n\neq m$, we have}
		d(x_n,x_m)&\geq\eps.\end{align*}
		By (\ref{conv_subseq}), $(x_n)_\nin$ contains a convergent subsequence $(x_{n_k})_{k\in\N}$. 
		Let $x:=\lim\limits_{k\to\infty}x_{n_k}$. Then $d(x_{n_k},x)<\frac\eps2$ for all $k>k_0$, hence $d(x_{n_k},x_{n_l})<\eps$
		for all $k,l>k_0$ $\lightning.$
\end{proof}

\begin{lemma}\label{sub_bounded_compact} Let $(X,d)$ be a metric space, and let $A\sub X$, $A\neq\emptyset$.
	\begin{enumerate}[(i)]
	\item\label{subbounded} If $X$ is totally bounded, then also $A$ is totally bounded.
	\item\label{closbounded} If $A$ is totally bounded, then also $\overline A$ is totally bounded.
	\end{enumerate}
\end{lemma}
\begin{proof}~
(\ref{subbounded}). Let $\eps>0$. By hypothesis, there exists an $\frac\eps2$-net $\{x_1,\ldots x_n\}$ of $X$.
			WLOG, let $A\cap U_{\frac\eps2}(x_j)\neq\emptyset$ if and only if $1\leq j\leq m$, $m\leq n$. For each $1\leq j\leq m$, choose
			$y_j\in A\cap U_\frac\eps2(x_j)$. Let $y\in A$. Then there exists $1\leq j\leq m$ with $y\in U_\frac\eps2(x_j)$, and hence
			\begin{align*}d(y,y_j)\leq d(y,x_j)+d(x_j,y_j)<\eps.\end{align*}
			This implies that $\{y_1,\ldots,y_n\}$ is an $\eps-$net for $A$.
			
(\ref{closbounded}). Let $\eps>0$. By hypothesis, there exists an $\frac\eps2-$net $\{y_1,\ldots,y_n\}$ for $A$.
			Let $x\in\overline A$. Then there exists $y\in A$ with $d(x,y)<\frac\eps2$. Let $y_j$ be such that $d(y,y_j)<\frac\eps2$.
			This yields
			\begin{align*}d(x,y_j)\leq d(x,y)+d(y,y_j)<\eps,\end{align*}
			hence $\{y_1,\ldots,y_n\}$ is an $\eps-$net for $\overline A$.
\end{proof}

\begin{kor}\label{cor_comp_bounded}
	Let $(X,d)$ be a complete metric space, and let $A\sub X$. Then the following are equivalent.
	\begin{enumerate}[(i)]
	\item\label{closcompact} $\overline A$ is compact.
	\item\label{tbounded} $A$ is totally bounded.
	\end{enumerate}
\end{kor}
\begin{proof}~
(\ref{closcompact})$\Ra$(\ref{tbounded}). Since $\overline A$ is compact, by \ref{bounded_compact}, $\overline A$ is
		totally bounded. By Lemma \ref{sub_bounded_compact}, $A$ is totally bounded.
		
(\ref{tbounded})$\Ra$(\ref{closcompact}). Since $A$ is totally bounded, by Lemma \ref{sub_bounded_compact},
		$\overline A$ is totally bounded. Since $X$ is compact, $\overline A$ is also complete. Hence Theorem \ref{bounded_compact}
		implies that $\overline A$ is compact.
\end{proof}

\begin{definition}Let $(X,d)$ and $(X',d')$ be metric spaces, and let $f\colon X\to X'$.
	\begin{enumerate}[(1)]
	\item $f$ is \emph{continuous} in $x\in X$, if for each $\eps>0$ there exists $\delta >0$ such that 
		$d(x,y)<\delta\Ra d'(f(x),f(y))<\eps$ for all $y\in X$.
	\item $f$ is a \emph{homeomorphism}, if $f$ is bijective and $f$ and $f^{-1}$ are both continuous. $f$ is an \emph{isometry}, if
		$f$ is bijective and $d(x,y)=d(f(x),f(y))$ for all $x,y\in X$.
	\item $f$ is \emph{uniformly continuous}, if for each $\eps>0$ there exists $\delta>0$ with
		$d(x,y)<\delta\Ra d'(f(x),f(y))<\eps$ for all $x,y\in X$.
	\end{enumerate}
\end{definition}

Two metric spaces being isometric is a strong notion of equivalence for metric spaces, being homeomorphic is the properly weaker topological equivalence of metric spaces.

\begin{lemma} Let $(X,d)$ and $(X',d')$ be metric spaces, and let $f\colon X\to X'$.
	\begin{enumerate}[(i)]
	\item $f$ is continous $\eq$ $f^{-1}(U)$ is open in $X$ for all $U\sub X'$ open $\eq$ $f(x_n)\to f(x)$ for all $x_n\to x$ in $X$.
	\item Let $X$ be compact and $f$ continuous. Then $f$ is automatically uniformly continuous.
	\end{enumerate}
	\label{lem:openfuncarecont}
\end{lemma}
\begin{proof} Exercises.\end{proof}

We want to relate the relative compactness, that is the compactness of the closure, of a set of continuous real functions to the pointwise relative compactness of these functions. The relatively compact sets in $\R$ are the bounded sets by Heine-Borel theorem. If a set of continuous real functions is relatively compact, we obtain pointwise relative compactness, by continuity of $C(X) \to \R, f \mapsto f(x)$ for every $x\in X$. However, to prove the converse a second condition is needed: the equicontinuity of the functions.

\begin{definition}
	$F\sub C(X)$ is \emph{equicontinuous in $x\in X$}, if for each $\eps>0$ there exists a neighbourhood $U$ of $x$ with $|f(x)-f(y)|<\eps$
	for all $y\in U$ and $f\in F$. $F$ is called \emph{equicontinuous}, if it is equicontinuous in each $x\in X$.
\end{definition}

\begin{theo}[Arzela-Ascoli]\label{arzela-ascoli}
	Let $X$ be a compact metric space and $F\sub C(X)$. Then the following are equivalent.
	\begin{enumerate}[(i)] 
	\item\label{clos_compact} $\overline F$ is compact.
	\item\label{equicont} $F$ is equicontinuous and pointwise bounded.
	\end{enumerate}
\end{theo}
\begin{proof}
(\ref{clos_compact})$\Ra$(\ref{equicont}). Exercise.

(\ref{equicont})$\Ra$(\ref{clos_compact}). Let $F$ be equicontinuous and $F(x)\in\K$ bounded for all $x\in X$. Since $C(X)$ is complete, by \ref{cor_comp_bounded} it remains to prove that $F$ is totally bounded. For this, let $\eps>0$, and, for each $x\in X$, let $U_x$ be an open neighbourhood of $x$ with
\begin{align*}
|f(y)-f(x)|<\frac\eps2~\text{for all }f\in F\text{ and }y\in U_x.
\end{align*}
Let now $x_1,\ldots,x_n\in X$ be chosen such that $X=\bigcup\limits_{i=1}^n U_{x_i}$ and set
\begin{align*}
K&:=\bigcup\limits_{i=1}^n f(x_i)\sub\K.
\intertext{Since $K$ is bounded, there exists $\lambda_1,\ldots,\lambda_n\in\K$ with}
K&\sub\bigcup\limits_{j=1}^n U_\frac\eps2(\lambda_j).
\end{align*}
Define $\Phi$ to be the set of maps $\ph\colon\{1,\ldots,n\}\to\{1,\ldots,m\}$. Also, for $\ph\in\Phi$, set
\begin{align*}
F_\ph:=\{f\in F:|f(x_i)-\lambda_{\ph(i)}|<\frac\eps3~\text{for }1\leq i\leq n\}.
\end{align*}
Then
\begin{align*}
F=\bigcup\limits_{\ph\in\Phi}F_\ph.
\end{align*}
To see this, note that for $f\in F$ and each $1\leq i\leq n$, there exists $\ph(i)\in\{1,\ldots,m\}$ with $f(x_i)\in U_\frac\eps6(\lambda_{\ph(i)})$. Hence $f\in F_\ph$. For $f,g\in F_\ph$ and $y\in U_{x_i}$, we then obtain
\begin{align*}
|f(y)-g(y)|\leq|f(y)-f(x_i)|+|f(x_i)-\lambda_{\ph(i)}|+|\lambda_{\ph(i)}-g(y)|\leq\eps.
\end{align*}
Thus $d(f,g)\leq\eps$ for all $f,g\in F_\ph$, and hence a finite $\eps-$net does exist.
\end{proof}
	

\section{Normed Spaces}

\begin{definition} Let $E$ be a linear space over $\K$.
	\begin{description}
	\item{(1)} Then a map $\|\cdot\|\colon E\to[0,\infty)$ is called a \emph{norm an $E$}, and $(E,\|\cdot\|)$ a \emph{normed space}, 
		if for all $x,y\in E$, $\lambda\in\K$
		\begin{description}
		\item{(i)} $\|x\|=0\eq x=0$.
		\item{(ii)} $\|\lambda x\|=|\lambda|\cdot\|x\|.$
		\item{(iii)} $\|x+y\|\leq\|x\|+\|y\|.$
		\end{description}
		$E$ is called a \emph{Banach space}, if $(E,d_{\|\cdot\|})$ is complete.
	\item{(2)} Two norms $\|\cdot\|_1$ and $\|\cdot\|_2$ are \emph{equivalent} if there exist $\alpha,\beta >0$ such that
		\begin{align*} \alpha \|x\|_1\leq\|x\|_2\leq\beta\|x\|_1\text{ for all }x\in E.\end{align*}
	\end{description}
\end{definition}

\begin{bem}
	Let $E$ be a linear space.
	\begin{enumerate}[(1)]
		\item If $\| \cdot \|$ is a norm on $E$, then $d_{\|\cdot\|}(x,y) = \|x-y\|$ defines a metric on $E$.
		\item Let $\| \cdot \|_{1}$ and $\| \cdot \|_{2}$ be equivalent. Then $(E,\| \cdot \|_{1})$ is complete
			if and only if $(E,\| \cdot \|_{2})$ is complete.
		\item $ | \|x\| - \|y\| | \le \|x - y\|$. In particular $\| \cdot \|: E \mapsto \R$ is Lipschitz-continous.
		\item The algebraic operations
			\begin{itemize}
				\item $+: E \times E \mapsto E \qquad (x,y) \mapsto x + y$
				\item $\cdot: \K \times E \mapsto E \qquad (\lambda,y) \mapsto \lambda \cdot y$
			\end{itemize}
			are continous:
			\begin{itemize}
				\item $\|(x+y) - (x_0 + y_0)\| \le \| x-x_0 \| + \|y-y_0\|$
				\item $\| \lambda x - \lambda_0 x_0 \| \le |\lambda| \|x - x_0\| + |\lambda - \lambda_0| \|x_0\| $
			\end{itemize}
		\item If $F \subset E$ is a subspace, so is $\overline F$
	\end{enumerate}
\end{bem}


% This has no numeration in the skript. So the numeration does not fit starting from here.
\begin{definition}[Quotientspace]
	Let $E$ be a linear space, $F \subset E$. Then
	\begin{equation*}
		x \sim y :\eq x-y \in F \qquad (x,y \in E)
	\end{equation*}
	is an equivalence relation.
	\begin{equation*}
		[x]_{\sim} = \left\{ y \in E: y-x \in F \right\} = \left\{ y \in E : y \in x + F\right\} = x + F
	\end{equation*}
	So $[x]_{\sim}$ is an \emph{affine subspace}.

	The \emph{quotient space} $E / F$ is defined by
	\begin{equation*}
		E/F := \left\{ x+F : x \in E \right\}
	\end{equation*}
	via
	\begin{equation*}
		[x]_{\sim} + [y]_{\sim} := [x+y]_{\sim} \qquad (x + F) + (y + F) := ( (x+y) + F )
	\end{equation*}
	and
	\begin{equation*}
		\lambda [x]_{\sim} := [\lambda x]_{\sim} \qquad \lambda (x + F) := ( (\lambda x) + F )
	\end{equation*}
	the space $E/F$ becomes a linear space.
\end{definition}

\begin{lemma}
	Let $(E, \|\cdot\|)$ be a normed space and let $F \subset E$ be a \emph{closed} subspace. Then
	\begin{equation*}
		\|x+F\| := \inf\{ \|x+y\| : y \in F\}
	\end{equation*}
	defines a norm on $E/F$. Moreover, if $E$ is Banach space so is $E/F$.
\end{lemma}

\begin{proof}
	\begin{enumerate}[(i)]
		\item Let $\|x + F\| = 0$, this implies it exists $\sequ{y} \subset F$, such that 
			\begin{equation*}
				\| x - y_n \| \xrightarrow{n \to \infty} 0
			\end{equation*}
			Since $y_n \in F$, $F$ closed and $x \in F$
			\begin{equation*}
				x + F = F + F = F = [0]_{\sim} = 0 + F = 0
			\end{equation*}
		\item $\|\lambda (x + F)\| = \| (\lambda x ) + F\| = \inf\{\|\lambda x + y : y \in F\|\}$

			For $\lambda = 0$ we have:
			\begin{equation*}
				\| \lambda (x + F) \| = 0 = |\lambda| \|x + F \|
			\end{equation*}

			And for $\lambda \neq 0$:
			\begin{align*}
				\|\lambda (x+F)\|
				&= \inf\left\{ \|\lambda x + y\| : y \in F \right\}\\
				&= |\lambda| \inf\left\{ \| x + y\| : y \in F \right\}\\
				&= |\lambda| \|x + F\|
			\end{align*}
		\item Let $x,y \in E, \epsi$. Choose $z_1, z_2$ such that,
			\begin{align*}
				\| x + F \| &\ge \|x + z_1\| - \frac{\eps}{2} \\
				\| y + F \| &\ge \|y + z_2\| - \frac{\eps}{2} \\
				\intertext{which gives us}
				\|(x+F) (y+F)\|	&= \|(x+y) + F\| \\
								&\le \|x+z_1 + y + y_2\| \\
								&\le \|x + F\| + \| y + F\| + \eps
			\end{align*}
			Let $E$ be complete and let $(x_n + F)_{n\in \N}$ be a Cauchy-sequence in $E/F$, i.e.
			\begin{equation*}
				\forall {\epsi} \quad \exists {N \in \N} \quad \forall {n,m \ge N} : \|(x_n - x_m) + F\| \le \eps
			\end{equation*}
			So for all $i \in \N$ we can find $n_i$, such that:
			\begin{align*}
				\|x_{n_{i+i}} - x_{n_i} + F\| 		&\le 2^{-i} \\
				\intertext{in particulat it exists $y_i \in F$ such that}
				\|x_{n_{i+1}} - x_{n_{i}} + y_i\|	&\le 2^{i}
			\end{align*}
			We may assume $n_i < n_{i+1}$. Now define
			\begin{align*}
				z_1 	&:= 0 \\
				z_{i+1}	&:= y_i - z_i \qquad i \ge 1
			\end{align*}
			In particular we have $y_i = z_{i+1} - z_i$:
			\begin{equation*}
				\|(x_{n_{i+1}} + z_{i+1}) - (x_{n_i} + z_i)\| < 2^{-i}
			\end{equation*}
			Now we define $\eta_i := x_{n_i} + z_i$, which gives us
			\begin{align*}
							& \| \eta_{i+1} - \eta_i\| < 2^{-1}\\
				\Ra \quad	& \|\eta_{m+k} - \eta_m\| \le \sum_{i=0}^{k-1} \|\eta_{m+i+1} - \eta_{m+i}\| < \sum_{i=0}^{k+1} 2^{-m-1} \le 2^{1-m}\\ 
				\Ra \quad	& \sequ{\eta} \text{ is a Cauchy-sequence in } E\\
				\Ra \quad	& \sequ{\eta} \text{ converges}
			\end{align*}
			Now we set $\lim_{n \to \infty} \eta_n =: x$.
			We obtain:
			\begin{align*}
											&\|(x_n + F) - (x + F)\| \\
											=& \|(x_n - x) + F\| \\
				\stackrel{z_i \in F}{\le}	& \| x_{n_i} + z_i - x\|\\
											=& \| \eta_i - x\| \to 0
			\end{align*}
			Which gives us a convergent subsequence, so the Cauchy-sequence is covergent itself.
	\end{enumerate}
\end{proof}

\begin{lemma}\label{lem:banachspaceIfQuotientSpacesAre}
	Let $E$ be a normed space, $F \subset E$ closed subspace. If $F$ and $E/F$ are Banach spaces
	so is $E$.
\end{lemma}

\begin{proof}
	Let $\sequ{x} \subset E$ be a Cauchy-sequence in $E$. Which gives us.
	\begin{align*}
		\|(x_n + F) - (x_m + F)\| = \|(x_n + x_M) + F \| \le \|x_n + x_m\|.
	\end{align*}
	So $(x_n + F)_n \subset E/F$ is a Cauchy-sequence in $E/F$. With $x+F := \lim_{n \to \infty} x_n + F$.
	we obtain:
	\begin{align*}
		\Ra& \inf\left\{ \|x_n - x + y\| : x \in F \right\} = \|(x_n - x) + F\| \to 0\\
		\Ra& \exists \sequ{y} \subset F : \|x_n - x + y_n\| \to 0\\
		\Ra& \|y_n - y_m\| &= \|y_n + x_n - x - x_n + x_m - y_m - x_m + x\|\\
		&&\|y_n + x_n - x\| + \|x_n - x_m\| + \|y_m + x_m - x\| \xrightarrow{m,n \to \infty} 0\\
		\Ra& \sequ{y} \text{ is a Cauchy-sequence in } F.\\
		\Ra& y := \lim_{n \to \infty} y_n \in F \text{ exists}\\
		\Ra& \|x_n - x + y\| \le \|x_n + y_n - x \| + \|y - y_n\| \xrightarrow{n \to \infty} 0\\
		\Ra& x_n \to (x-y), n \to \infty
	\end{align*}
\end{proof}

\begin{kor}
	A finite-dimensional normed space $E$ is always a Banach-space.
\end{kor}

\begin{proof}
	Proof by induction.
	\begin{description}
		\item[$n=1$] Let $\dim{E} = 1$. Choose $x \in E$ such that $\|x\|=1$.
			Then $q: \R \mapsto E, q(\lambda) := \lambda x$ is isometric. So $E \sim \R$
			and $E$ is a Banach-space.
		\item[$n\to n+1$] Let $\dim{E} = n+1$. Choose $x \in E \setminus \{0\}$ and set $F := \Span{x}$.
			Because $\dim{F} = 1$ we know that $F$ is complete, so closed. Now consider:
			\begin{equation*}
				\dim(E / F) = \dim E - \dim F = n
			\end{equation*}
			By assumption we get that $E / F$ is complete and by lemma \ref{lem:banachspaceIfQuotientSpacesAre}
			we see that $E$ is a Banach space.
	\end{description}
\end{proof}

\begin{lemma}
	Let $F$ be a closed subspace of a normed space $E$. Then for each $x \in E \setminus F$
	there exists $M,M'$ such that,
	\begin{align*}
		\forall_{y\in F, \lambda \in \K} : |\lambda| &\le M |\ \lambda x + y \| \\
										\|y\| &\le M'\ \|\lambda x +y\|
	\end{align*}
\end{lemma}

\begin{proof}
	For $x \nin F$ we have $\|x + F\| \neq 0$. We set
	\begin{align*}
		M	&:= \|x + F\|^{-1} \\
		M'	&:= 1 + M \|x\|
	\end{align*}
	Then for $y \in F$ and $\lambda \in \K$ we have
	\begin{align*}
		\|\lambda\|	&\le	M |\lambda| \|x + F\| \\
					&=		M \|\lambda x + y\| \\
					&\le	M \|\lambda x + y\| \\
		\intertext{For $\lambda \neq 0$ we get}
		\|y\|		&\le	\|y + \lambda x\| + |\lambda| \|x\| \\
					&\le	\|y + \lambda x\| + M \|\lambda x + y\| \|x\| \\
					&=		\|y + \lambda x\| ( 1 + M \|x\|)
	\end{align*}
\end{proof}

\section{Linear Operators, Dual Space}

\begin{lemma}
	Let $E,F$ be normed spaces over $\K$ and $T: E \rightarrow F$ be a linear operator (a linear map). Then the following are 			equivalent: \label{bounded_continous}

	\begin{enumerate}[(i)]
		\item $T$ is continuous on $E$.\label{cond:continous}
		\item $T$ is continuous in one point $x_0 \in E$. \label{contpoint}
		\item $T$ is bounded i.e. $\norm{Tx} \le c\norm{x}$ for all $x \in E$ and some $c \in \R_{\ge 0}$. \label{cond:bounded}
	\end{enumerate}
\end{lemma}

\begin{proof}
	\eqref{cond:continous} $\Ra$ \eqref{contpoint}: This is trivial.
	
	\eqref{contpoint} $\Ra$ \eqref{cond:bounded}: Let $T$ be continuous in $x_0$. Then there exists $\delta \in \R_{> 0}$ such that $\norm{x - x_0} \le \delta \Ra \norm{Tx - Tx_0} \le 1$ for all $x \in E$. Let $y = \delta^{-1}\cdot(x_0 - x)$ then $\norm{\dfrac{\delta \cdot y}{\norm{y}}} \le \delta$. We obtain $\norm{Tx_0 - Tx} = \norm{T(\dfrac{\delta \cdot y}{\norm{y}})} \le 1. \iff \norm{Ty} \le \dfrac{\norm{y}}{\delta}$.
	
	\eqref{contpoint} $\Ra$ \eqref{cond:continous}:  Let $\varepsilon > 0$ then $\norm{Tx} \le c \norm{x} < \varepsilon \La \norm{x} \le \dfrac{\varepsilon}{c}$. This implies $T(K_{\dfrac{\varepsilon}{c}}(x)) \subseteq K_{\dfrac{\varepsilon}{c}}(Tx)$.

\end{proof}

\begin{definition}
Let $E,F$ be normed spaces over $\K$. Then $\L(E,F) = {T:E \rightarrow F \ | T\,is \,linear\,and\,bounded}$. If $E = F$ we write $\L(E)$. We define the operator norm on $\L(E,f)$ by $\norm{T}=sup_{\norm{x}_{E} \le 1} \norm{Tx}_{F}$.
\end{definition}

\begin{lemma}
Let $E,F,G$ be normed spaces over $\K$:
\begin{enumerate}
	\item[(i)] $(\L(E,F), \norm{\cdot})$ is a normed linear space. Also for $T \in \L(E,F)$ we have $\norm{T}= \sup_{\norm{x}=1} \norm{Tx} = \sup_{x \neq 0} \dfrac{\norm{Tx}}{\norm{x}}$.
	\item[(ii)] If $T \in \L(E,F)$ and $S \in \L(F,G)$ then $T \circ S \in \L(E,G)$ and $\norm{S \circ T} \le \norm{S} \norm{T}$.
\end{enumerate}
\end{lemma}

\begin{proof}
Excercise.
\end{proof}

\begin{definition}
Let $E,F$ be normed spaces over $\K$. Then the space of continuous linear functions $\L(E,\K)$  is the dual space $E^{*}$ of $E$.
\end{definition}

\begin{theo}
Let $E$ be a normed space over $\K$ and let $F$ be a Banach space over $\K$. Then $\L(E,F)$ is a Banach space. In particular $E^{*}$ is a Banach space.
\end{theo}

\begin{proof}
Let $\sequence{T_n}{n}$ be a Cauchy sequence in $\L(E,F)$. Then since $\norm{T_n x - T_m x} \le \norm{T_n - T_m}\cdot \norm{x}$, also $\sequence{T_n x}{n}$ is a Cauchy sequence in $F$ for each $x \in E$. Since $F$ is complete we can find $Tx = \lim_{n \rightarrow \infty} T_n x$. This defines $T:E \rightarrow F$.
\begin{enumerate}
	\item[1)]Let $x,y \in E$ $\lambda, \mu \in \K$ then $T(\lambda x + \mu y) = \lim_{n \rightarrow \infty} T_n(\lambda x + \mu y) =  \lambda \cdot \lim_{n \rightarrow \infty} T_n x + \mu \cdot \lim_{n \rightarrow \infty} T_n y = \lambda \cdot T x + \mu \cdot T y $
	\item[2)] Let $\varepsilon \in \R_{>0}$ $N_{\varepsilon} \in \N$ such that $\norm{T_n - T_m} < \varepsilon \, \forall m,n \ge N_{\varepsilon}$. This implies $\norm{T_n x - T_m x} \le \norm{T_n - T_m} \norm{x} < \varepsilon \norm{x} \, \forall m,n \ge N_{\varepsilon}$ and $\forall x \in E$. Letting $n \rightarrow \infty$ we obtain $\norm{Tx - T_{m} x} \le \varepsilon \cdot \norm{x} \forall m \ge N_{\varepsilon}$ hence $\norm{Tx - T_{N_{\varepsilon} x}} \le \varepsilon \cdot \norm{x}$. This implies $\norm{Tx} \le \varepsilon \cdot \norm{x} + \norm{T_{ N_{\varepsilon}} x} \le (\varepsilon + \norm{T_{N_{\varepsilon}}}) \cdot \norm{x}$ $\forall x \in E$. Hence $T \in \L(E,F)$
	\item[3)] $\norm{Tx - T_{N_{\varepsilon} x}} \le \varepsilon \cdot \norm{x} \forall m \ge N_{\varepsilon}$ $\norm{T - T_{N_{\varepsilon}}} \le \varepsilon$. Hence  $\lim_{m \rightarrow \infty} T_m = T$.
\end{enumerate}
\end{proof}

\begin{theo}
Let $E$ be a normed space over $\K$ and let $F$ be a Banach space over $\K$, $L \subseteq E$ a linear subspace, $F$ be a Banach space, and $T: L \rightarrow F$ a continuous linear operator. Then theere exists $S: \bar{L} \rightarrow F$ with $S_{| \, L} = T$, $S$ is continuous, linear and  $\norm{S} = \norm{T} $, $S \in \L(\bar{L}, F)$.
\end{theo}

\begin{proof}
$ \ $
\begin{enumerate}
	\item[1)]Let $x \in \bar{L}$ and  $\sequence{x_n}{n} \subseteq L$, $x = \lim_{m \rightarrow \infty} x_n $ we observe: $\norm{Tx_n - Tx_m} \le \norm{T} \norm{x_n - x_m}$. This implies that $\sequence{Tx_n}{n}$ is a Cauchy sequence. Since $F$ is a Banach space, hence $\sequence{Tx_n}{n}$ converges. Now chosse $\sequence{y_n}{n} \subseteq L$ with $x= \lim_{n \rightarrow \infty} y_n$ then $\norm{Tx_n - Ty_n} \le \norm{T} \norm{x_n - y_n}$, which implies $\lim_{n \rightarrow \infty} Tx_n = \lim_{n \rightarrow \infty} Ty_n$. This means tahat the map $S$ defined by $Sx = \lim_{n \rightarrow \infty} Ty_n$ is well defined.
	\item[2)] Let $x,y \in \bar{L}$ be arbitrary, choose $\sequence{x_n}{n}, \sequence{y_n}{n} \subseteq L$ with $\lim_{n \rightarrow \infty} x_n = x$, $\lim_{n \rightarrow \infty} y_n = y$ and  $\lambda , \mu \in \K$. Therefore $S(\lambda x + \mu y) = \lim_{n \rightarrow \infty} T(\lambda x_n + \mu y_n) = \lambda \lim_{n \rightarrow \infty} Tx_n + \mu \lim_{n \rightarrow \infty} Ty_n = \lambda Sx + \mu Sy$ , which implies that $S$ is linear.
	\item[3)] Let $x \in \bar{L}$, $\lim_{n \rightarrow \infty} x_n = x$, $x_n \in L$, $\norm{x} \le 1$. Then $\norm{Sx} = norm{\lim_{n \rightarrow \infty} Tx_n} = \lim_{n \rightarrow \infty} \norm{Tx_n}=\norm{Tx} \le \norm{T}$. Therefore $\norm{S} \le \norm{T}$. Since $L \subseteq \bar{L}$ $S_{| \, L} = T$ we have $\norm{T} \le \norm{S}$ hence $\norm{S} = \norm{T}$.
	\item[4)] From 3) it follows that $\norm{Sx} \le \norm{S} \norm{x} = \norm{T} \norm{x} < \infty$. This implies that $S$ is contiuous.
	\item[5)] Now choose $R \in \L(\bar(L), F)$ with $R_{| \, L} = T$. Then all $x \in \bar{L}$ and sequences $\sequence{x_n}{n}$ with $\lambda \lim_{n \rightarrow \infty} x_n = x$ satisfy $Rx = \lim_{n \rightarrow \infty} Rx_n = \lim_{n \rightarrow \infty} Tx_n = S_x$, which shows the uniqueneness of $S$.
\end{enumerate}
\end{proof}

\begin{lemma} \label{lift_from_dense}
    Let $E$ be a normed space over $\mathbb{K}$, $L \in E$ a subspace, $F$ a   	
    Banach space and $T: L \to F$ a continous linear operator. Then there exists a unique $S \in \scrL(\overline{L},F)$    
    with 
    $S_{\vert_L}=T$. There holds
		\begin{align*}
			\Vert S \Vert = \Vert T \Vert
		\end{align*}
\end{lemma}

\begin{proof}
	Let $x \in \bar{L}$ and let $(x_n) \subseteq L$ be a sequence which converges to x. We observe
	\[
      \Vert Tx_n - Tx_m \Vert \leq \Vert T \Vert \Vert x_n - x_m \Vert
     \] 
	This implies that $(Tx_n)$ is a Cauchy sequence in $F$. $F$ is complete - hence $(Tx_n)$ converges.
	For any other sequence $(y_n)$ which also converges to $x$ we have
	\[
		\Vert Tx_n - Ty_n \Vert \leq \Vert T \Vert \Vert x_n - y_n \Vert
	\]
	Thus $\lim Tx_n = \lim Ty_n$. We can now define
    \[
		Sx = \lim Tx_n
	\]
	and have no concernes about well-defining issues. 

	It follows immediately $S_{\vert_L}=T$, because for $x \in L$ we can just choose the constant sequence $x_n:=x , \forall n$ as 	''defining sequence''. The linearity of $S$ is also easily proven:
	\[
		S(x+y) = \lim T(x_n+y_n) = \lim Tx_n + Ty_n =Sx +Sy
	\]
	and analogously $S(\lambda x) = \lambda Sx$.

	Moreover, $\Vert S \Vert \geq \Vert S_{\vert_L} \Vert = \Vert T \Vert$. To prove the reverse inequality, let $x \neq 0 \in 			\bar{L}$ and $(x_n)$ be a sequence in $L$ converging to $x$. For large $n$, $x_n \neq 0$ and therefore we can consider
	\[
		\frac{\Vert Sx \Vert}{\Vert x \Vert} = \lim \frac{\Vert\ T x_n\Vert }{\Vert x_n \Vert } 
												\leq \lim \Vert T \Vert = \Vert T \Vert
	\]
	This proves $\Vert S \Vert = \Vert T \Vert$, thus $S$ is bounded.

	It only remains to prove the uniqueness of $S$. Consider another continous linear operator $R$ with 
	$R_{\vert_L}=T$, and $x \in \bar{L}$. 
	For any sequence converging to $x$ there follows
	\[
		Rx= \lim Rx_n = \lim Tx_n =Sx
	\]
	We used the continuity of $R$.
\end{proof}

We draw a corollary. It follows from the uniqueness part of the lemma.

\begin{kor} \label{cor_lift_from_dense}
	If two bounded linear operators $S,T \in \scrL(E,F)$, where $F$ is a Banach space, 
	coincides in a dense subspace of $E$, then they coincide in $E$
\end{kor}

\begin{lemma} \label{inverse}
	Let $E,F$ be normed spaces over $\mathbb{K}$ and $T: E \to F$ linear. Then the following are equivalent
	\begin{enumerate}[(i)]
		 \item There exists a linear, cond:continousinous, inverse operator \label{invertible}
 			\[
  				T^{-1}: T(E) \to E
 			\]
 		\item  There exists $c >0$ so that $c \Vert x \Vert \leq \Vert Tx\Vert$ \label{inv_bound}
	\end{enumerate}
\end{lemma}

\begin{proof}
 	 \eqref{invertible} $\Ra$ \eqref{inv_bound} Assume $T^{-1}$ exists. The cond:continousinuity of it gives us the 
 	 existence of a $\gamma>0$ such that:
 	\[
 		\Vert T^{-1}y \Vert \leq \gamma \Vert y \Vert
 	\]
	for an arbitrary $x$ we put $y=Tx$ to obtain
	\[
		\Vert x \Vert \leq \gamma \Vert Tx\Vert
	\]
	Putting $c := \frac{1}{\gamma}$ we have proven \eqref{inv_bound}.
	
	\eqref{inv_bound} $\Ra$ \eqref{invertible} We observe that \ref{invertible} secures the injectivity of $T$ (if $ x \in \ker(T)$, 		then $\Vert x \Vert =0$). Thus $T^{-1}: \ran E \to E$ exists.

	Now letting $y=Tx$ in \eqref{invertible} assures the existence of a $c>0$ with
	\[
		c\Vert T^{-1}y \Vert \leq \Vert y \Vert \quad \forall y \in T(E)
	\]
	Thus the inverse operator is cond:continousinous.
\end{proof}

\begin{definition}[Graph of $T$] \label{graph}
	Let $E,F$ be normed spaces, $L \subseteq E$ a subspace and $T: L \to F$ a linear operator. 
	\begin{enumerate}[(i)]
	\item
		\[
			G_T = \{(x,Tx) , x\in L \} \subseteq L \times F
		\]
		is called the \emph{graph of $T$}.
	
	\item If $G_T$ is closed, $T$ is \emph{closed}.
	\end{enumerate}
\end{definition}

\begin{lemma}
	Let $E,F,T$ and $L$ be as above. Then the following are equivalent
	\begin{enumerate}[(i)]
		\item{$T$ is closed \label{closed_graph}}
		\item{ If $(x_n) \subseteq(L)$ converges to $x \in E$ and $(Tx_n)$ to $y \in F$, then $x \in L, y=Tx$ \label{seqcrit}}
	\end{enumerate}
\end{lemma}

\begin{proof}
	Since
	\[
		\Vert (x_n, Tx_n) - (x,y) \Vert = \max(\Vert x_n - x \Vert, \Vert Tx_n -y \Vert )
	\]
	we have that if $x_n \to x, Tx_n \to y$, then
	\[
		\lim_{n \to \infty} (x_n, Tx_n) = (x,y)
	\]
	Because of $G_T$ being closed $(x,y) \in G_T$. Thus $x \in L, y=Tx$.

	Now consider a convergent sequence $(x_n, y_n) \to (x,y)$ in $G_T$. Because of the convergence of $y_n =Tx_n$ and  	
	\eqref{seqcrit}, there follows $x \in L, y=Tx$, thus $(x,y) \in G_T$.
\end{proof}

\begin{rem}
	If $L$ is closed, and $T$ is cond:continousinous, then $T$ is closed. In particular, all $T \in \scrL(E,F)$ are closed.
\end{rem}
\begin{proof}
	If $(x_n,Tx_n) \to (x,y)$, then ($T$ closed) $x \in L$. Continuity of $T$ now implies $Tx_n \to Tx$, thus $(x,y) \in G_T$.
\end{proof}

\begin{theo}
 	Let $1 \leq p < \infty$. Define $q$ so that $\frac{1}{p} + \frac{1}{q}=1$ , i.e.
 	\[
 	q= \begin{cases} \frac{p}{p-1} &: 1<p<\infty \\ \infty &: p=1 \end{cases}
 	\]
 	Moreover define for $y \in l_q$.
	\[
		f_y : l_p \to \mathbb{K}, x= (x_n) \mapsto \sum_{n=1}^{\infty} x_ny_n
	\]
	Then $f_y \in l_p^{\ast}$ and $y \mapsto f_y$ is an isomorphism.
	In particular $l_q \sim l_p^{\ast}$.
\end{theo}

\begin{proof}
	First of all, $f_y$ is welldefined (i.e. the series converges) because of the H"older-inequality:
    \[
        \sum_{n=1}^{\infty}\vert x_n y_n \vert \leq \Vert x \Vert_p \Vert y \Vert_q
    \]
    That $f_y$ is linear is evident. Furthermore, by the above, we have
	\[
		\Vert f_y(x) \Vert  \leq \Vert x \Vert_p \Vert y \Vert_q
	\]
	Thus $f_y$ bounded. We conclude that $f_y \in l_p^{\ast}$.

	We now claim that $\Vert f_y \Vert = \Vert y \Vert_q$. We already proved $\Vert f_y \Vert \leq \Vert y \Vert_q$. 
	To prove the inverse inequality, we consider the cases $p=1, p>1$ separately.

	\underline{$p=1$: } Let $\epsilon>0$. There exists an $n \in \mathbb{N}$ such that
	\[
		\vert y_n \vert \geq \Vert y \Vert_q -\epsilon
	\]
	Set $x=e_n \in l_1$. We obtain $f_y(x)=y_n$. Because of $\Vert x \Vert =1$ 
	there follows $\Vert f_y \Vert\geq \Vert y\Vert_{\infty} - \epsilon$. 
	Because of $\epsilon$ arbitrary, we conclude $\Vert f_y \Vert \geq \Vert y \Vert_q$

	\underline{$p>1$} Define $x$ by
	\[
	x_n = \begin{cases} 0 & \quad y_n=0 \\
                  \frac{\vert y_n \vert^q}{y_n} &\text{else}
			\end{cases}
	\]
	Then $\sum_{n=1}^{\infty}\vert x_n \vert^p = \sum_{n=1}^{\infty} \vert y_n \vert^{p(q-1)}=\sum_{n=1}^{\infty} \vert y_n 			\vert^{q}<\infty$, thus $x \in l_p$.
	
	We now compute $f_y(x)$
	\[
		f_y(x)=\sum_{n=1}^{\infty}x_ny_n= \sum_{n=1}^{\infty}\vert y_n \vert^q = \Vert y \Vert_q^q
	\]
	Thus $\frac{\vert f_y(x) \vert}{\Vert x \Vert} =\Vert y \Vert ^{\frac{q(p-1)}{p}}=\Vert y \Vert_q$, and we conclude $\Vert f_y 	\Vert \geq \Vert y \Vert_q$.

	This proves the injectivity of $y \mapsto f_y$. Now the surjectivity. Let $x \in l_p^{\ast}$ and put
	\[
		y_n:=f(e_n)
	\]
	
	To prove $y \in l_q$, we again treat the two cases $p=1$, $p>1$ separately.
	
	\underline{$p=1$} We have for all $n$
 	\[
		\vert y_n \vert = \vert f(e_n) \vert \leq \Vert f \Vert \Vert e_n \Vert = \Vert f\Vert
	\]
	Thus $\Vert y \Vert_{\infty} \leq \Vert f \Vert$, $y \in l_{\infty}$.
	
	\underline{$p>1$} For all $m \in \mathbb{N}$ there holds
	\[
		\sum_{n=1}^m \vert y_n \vert^q = \sum_{\substack{ n=1 \\ y_n \neq 0}}^m\frac{\vert y_n \vert^q}{y_n}f(e_n) = 
		f\Big(\sum_{\substack{ n=1 \\ y_n \neq 0}}^m\frac{\vert y_n \vert^q}{y_n}e_n \Big) 
		\leq \Vert f \Vert  \Vert\sum_{\substack{ n=1\\ y_n \neq 0}}^m\frac{\vert y_n \vert^q}{y_n}e_n\Vert_p
	\]
	We have 
	\[
		\Vert \sum_{\substack{ n=1 \\ y_n \neq 0}}^m\frac{\vert y_n \vert^q}{y_n}e_n \Vert_p = 
		\Big( \sum_{\substack{ n=1 \\ y_n \neq 0}}\vert y_n \vert^{p(q-1)} \Big)^{\frac{1}{p}}
	\]
	And thus
	\[
		\sum_{n=1}^m \vert y_n \vert^q \leq \Vert f \Vert \Big(\sum_{n=1}^m \vert y_n \vert^q \Big)^{ \frac{1}{p}}
	\]
	which implies
	\[
		\Big(  \sum_{n=1}^m \vert y_n \vert^q \Big)^{\frac{1}{q}} \leq \Vert f \Vert
	\]
	
	Letting $m \to \infty$, we get $\Vert y \Vert_q <\infty \Rightarrow y \in l_q$.
	
	Finally, $f=f_y$ because of for $x= \sum_{n=1}^mx_n e_n$, there holds
	\[
		f(x) = \sum_{n=1}^m x_n f(e_n) = \sum_{n=1}^m x_ny_n =f_y(x)
	\]
	Thus $f$ coincides with $f_n$ on the dense linear subspace $\text{span} (e_n)_{n\in \mathbb{N}}$, 
	hence (Lemma \eqref{lift_from_dense}, $\mathbb{K}$ is a Banach space) they are equal.
\end{proof}

\begin{lemma} \label{dual lemma}
	Let $E,F$ be normed spaces, and let $T \in \scrL(E,F)$- Then the operator $\dual{T}: \dual{F} \to \dual{E}$ defined by
		\[
		\dual{T}(\phi)(x)= \phi \circ T(x)
		\]
	satisfies $\dual{T} \in \scrL(\dual{F},\dual{E})$ and $\norm{\dual{T}}=\norm{T}$.
\end{lemma}

\begin{proof}
	$\dual{T}$ is obviously linear. Further
	\[
		\norm{(\dual{T}\phi)x} = \norm{\phi Tx \norm} \leq \norm{phi}\norm{T}\norm{x}
	\]
	This proves that $\dual{T}$ is bounded and $\norm{\dual{T}}\leq \norm{T}$. The inverse inequality is an exercize.
\end{proof}

\begin{def} \label{ dual def}
	Let $E,F$ be normed spaces and $T \in \scrL(E,F)$. The $\dual{T} : \dual{F} \to \dual{E} , \dual{T}\phi= \phi \circ T$ is 
	called the \emph{dual operator} of $T$
\end{def}

\section{Hahn-Banach Theorem and Corollaries}
 % To be overlooked %
\begin{definition}
\begin{enumerate}
	\item[a)] Let $E$ be a linear space. Then the algebraic dual space of $E$ which is the space of linear maps $E \rightarrow \K$, is denoted $E'$.
	\item[b)] Let $E$ be an $\R$-vector space. Then $p: E \rightarrow \R$ is a sublinear functional on $E$ if all $x,y \in E$ and $\lambda \in \R$ satisfy
	$p(x+y) \le p(x) + p(y)$ and $p(\lambda x) = \lambda p(x)$,
\end{enumerate}
\end{definition}

\begin{lemma}
... to be continued
\end{lemma}

% Axel Flinth ----------------------------
\paragraph*{WARNING} The first lemma of the day uses the Lemma of Zorn. If you're not comfortable with it, then essentially everything in this lecture will be very hard to swallow!

When now everyone has been warned it is time for

\begin{lemma} Let $E$ be an $\R$-vector space, $F \subseteq E$ a linear subspace and $p$ a sublinear functional on E. Further let $f \in F'$ with
\begin{equation*}
f(x) \leq p(x) \quad \forall x \in F
\end{equation*}
Then there exists $l \in E'$ with $l\vert_F =f$ and $l(x) \leq p(x) \forall x \in E$.
\end{lemma}
\begin{proof}
Set
\begin{eqnarray*}
\mathcal{L}= \{ (L,l)&:& L \text{ linear subspace of } E \text{ with } L\supset F  \\
                     && \text{and } l \in L' \text{ with } l\vert_ F =f \text{ and } l(x)\leq p(x) \forall x \in L\}
\end{eqnarray*}
To prove is the existence of a pair of the form $(E,l) \in \L$. To do this, we define an ordering on $\L$.
\begin{equation*}
(L_1, l_1) \leq (L_2,l_2) \Leftrightarrow L_1 \subset L_2 \land l_2\vert_{L_1} = l_1
\end{equation*}
We know that $\mathcal{L}\neq \varnothing$, then we have $(F,f) \in \L$. 

We want to use the Lemma of Zorn. Let $\mathscr{K}$ be a chain in $\L$. To prove is that $\KK$ has an upper bound in $\L$. We claim that if we set
\begin{equation*}
\tilde{L}= \{ L : \exists l' \text{ with } ( L,l') \in \mathscr{K} \}
\end{equation*} 
then $\tilde{L}$ is a linear subspace since $\KK$ is linearly ordered. define $\tilde{l}: \tilde{L} \to \R$ through
\begin{eqnarray*}
\tilde{l}(x) := l(x) &&\text{if } x \in L \text { and }\\
 && l \in L' \text{ with } (L,l) \in \mathscr{K}
\end{eqnarray*}
We know by the definition of $\tilde{L}$ that such an $l$ always exists. In fact, this definition is well-defined, since if $(L_2,l_2) \in \KK$ is another pair with $x \in L_2$, then one of the pairs is bigger with respect to the ordering ($\mathscr{K}$ is a chain), WLOG $(L_1,l_1) \leq (L_2,l_2)$. Then, because of $x \in L_1$ we have $l_2(x)=l_1(x)$.

$\tilde{l}\in \tilde{L}'$. To prove is only the linearity. Let $x_i \in L_i, \alpha_i \in \mathbb{R}, (L_i, l_i) \in \mathscr{K}, i=1,2$. Then WLOG $(L_1,l_1) \leq (L_2,l_2)$. Therefore $x_1  \in L_2$ and thus $\alpha_1 x_1 +\alpha_2 x_2 \in  L_2$. There follows
\begin{eqnarray*}
\tilde{l}(\alpha_1 x_1 + \alpha_2 x_2) &=& l_2(\alpha_2 x_1 +\alpha_2 x_2)  = \alpha_1 l_2(x_1) + \alpha_2 l_2(x_2) \\
 &=& \alpha_1 l_1(x_1) + \alpha_2 l_2(x_2) = \alpha_1 \tilde{l}(x_1) + \alpha_2 \tilde{l}(x_2)
\end{eqnarray*}


Finally $\tilde{l}(x) =l(x) \leq p(x)$. The Lemma of Zorn now provides the existence of a maximal element of $\mathcal{L}$ $(L,l) \in \mathcal{L}$. To prove is $L=E$. 

Suppose the opposite, then there exists $x_0 \in E \backslash L$. By Lemma 4.2, there exists $g \in (L + \R x_0)'$ with $g \vert_L =l$ and $g(x)\leq p(x) \quad \forall x \in (L+\mathbb{R}x_0)$. Then we although have $(L,l) < (L + \mathbb{R}x_0, g)$. Contradiction!

Thus $L=E$, this proves the Lemma.
\end{proof}

\begin{theo} Let $E$ be a vector space over $\mathbb{K}$, $F$ a linear subspace and $f \in F'$.

Let $p: E \to \mathbb{R}$ be a seminorm on E, i.e. for all $x,y \in E, \lambda \in \mathbb{K}$
\begin{eqnarray*}
p(x+y) &\leq & p(x) + p(y) \\
p(\lambda x) &=& \vert \lambda \vert p(x)
\end{eqnarray*}

Suppose that $\vert f(x) \vert \leq p(x)$ for all $x \in F$. Then there exists an $l \in E'$ with $l\vert_F=f$ and $\vert l(x) \vert \leq p(x) \quad \forall x \in E$.
\end{theo}

\begin{proof}
First consider $\mathbb{K}=\mathbb{R}$. Then $f(x) \leq p(x)$ for all $x \in F$ and $p(\alpha x )= \alpha p(x)$ for all $x\in E , \alpha \geq 0$. By Lemma 4.3, there exists some $l \in E'$ with $l\vert_F=f$ and $l(x)\leq p(x) \quad \forall x \in E$. Since also
\begin{equation*}
-l(x) = l(-x) \leq p(-x)=p(x)
\end{equation*}
we have $\vert l(x) \vert \leq p(x)$.

The case $\mathbb{K}=\mathbb{C}$ will be discussed in the exercises.
\end{proof}

\begin{theo}{(Hahn-Banach Theorem)} Let $E$ be a normed space, $F$ a linear subspace of $E$. Then for each $f \in F^{\ast}$ there exists some $l \in E^{\ast}$ with
\begin{equation*}
l \vert_F =f \land \Vert l \Vert = \Vert f \Vert
\end{equation*}
\end{theo}

\begin{proof}
Let $p$ be defined by
\begin{equation*}
p(x) = \Vert f \Vert \Vert x \Vert
\end{equation*}
Then $p$ is a seminorm - the properties are inherited from the norm properties of $\Vert \cdot \Vert$. Furthermore
\begin{equation*}
\vert f(x) \vert \leq p(x) \quad \forall x\in F
\end{equation*}
By Theorem 4.4, there exists $l \in E'$ with $l \vert_F=f$ and $ \vert l (x) \vert \leq p(x)= \Vert f \Vert \Vert x \Vert$. This proves in particular that $l \in E^{\ast}$ and $\Vert l \Vert \leq \Vert f \Vert$. Because of $l \vert_F=f$, the reverse inequality holds. Thus $\Vert l \Vert = \Vert f \Vert$. 
\end{proof}

\begin{kor}
Let $E$ be a normed space and $F$ a linear subspace of $E$ and $x \in E$ such that
\begin{equation*}
\delta := \inf_{y \in F} \Vert x - y \Vert >0
\end{equation*}
Then there exists an $l \in E^{\ast}$ with
\begin{equation*}
l \vert_F =0, \Vert l \Vert =1 \text{ and } l(x)=\delta
\end{equation*}
In particular, for any $x\neq0$ there exists an $l \in E^{\ast}$ with $\Vert l \Vert =1$ and $l(x)= \Vert x\Vert $.
\end{kor}
\begin{proof}
Let $G=F + \mathbb{K}x$ and $g: G \to \mathbb{K}$ be defined through
\begin{equation*}
g(y + \lambda x)= \lambda \delta \quad \forall y \in F, \lambda \in \mathbb{K}
\end{equation*}
$g$ is well defined, then because of $x \notin F$ we have $G = F \oplus \mathbb{K}x$. Further $g$ is linear, $g\vert_F=0$ and $g(x)=\delta$.

We know claim $\Vert g \Vert =1$. There holds
\begin{eqnarray*}
\vert g(y + \lambda x) \vert = \vert \lambda \vert \delta &=& \vert \lambda \vert \inf_{z \in F} \Vert z-x \Vert \\ &=& \inf_{z \in F} \Vert \lambda z - \lambda x \Vert = \inf_{z \in F} \Vert z+\lambda x \Vert \leq \Vert y + \lambda x \Vert
\end{eqnarray*}
thus $\Vert g \Vert \leq 1$.

Secondly, there exists for every $\epsilon>0$ a $z_{\epsilon} \in F$ with $\delta \leq \Vert z_{\epsilon} + x \Vert \leq \delta + \epsilon$. There follows
\begin{eqnarray*}
g(x + z_{\epsilon})&=& \delta \geq \Vert x+z_{\epsilon} \Vert - \epsilon \Leftarrow \\
g(\Vert x+ z_{\epsilon} \Vert^{-1}(x+z_{\epsilon})) &=&  \frac{\delta}{\Vert x+z_{\epsilon}\Vert} \geq 1 - \frac{\epsilon}{\Vert x + z_{\epsilon}} \geq 1-\frac{\epsilon}{\delta}
\end{eqnarray*}
Now apply Theorem 4.5 to lift $g$ up to $E^{\ast}$.

For the in-particular part, choose $F=\{0\}$.
\end{proof}

We know define an important concept
\begin{definition}
Let $E$ be a normed space, $M \subset E$ an arbitrary subset of $E$ and $L \in E^{\ast}$ one of $E^{\ast}$. Then the annihilator of $M$ in $E^{\ast}$ is
\begin{equation*}
M^{\perp} := \{ l \in E^{\ast} : l(x) =0 \forall x \in M \}
\end{equation*}
and the annihilator of $L$ in $E$ is
\begin{equation*}
L_{\perp} = \{ x \in E: l(x)=0 \forall l\in L \}
\end{equation*}
\end{definition}

\begin{rem} 
The annihilators are closed linear subspaces of $E^{\ast}$ and $E$, respectively. This follows from the continuity of $l \mapsto l(x), x \mapsto l(x)$.
\end{rem}

\begin{lemma} Let $E$ be a normed space and $\varnothing \neq M \subseteq E$. Then $(M^{\perp})_{\perp}$ is the closed linear hull of $M$, i.e. the smallest closed linear subspace of $E$ which contains $M$.
\end{lemma}
\begin{proof}
If $x \in M$, then $l(x)=0$ for all $l \in M^{\perp}$, thus $x \in (M^{\perp})_{\perp} \Leftarrow M \subseteq (M^{\perp})_{\perp}$.

Now let $F$ be the closed linear hull of $M$. By the remark, $F \subseteq (M^{\perp})_{\perp}$. Now assume there exists $x \in (M^{\perp})_{\perp} \backslash F$. Corollary 4.6 secures the existence of an $l \in (M^{\perp})_{\perp}^{\ast}$ with $l \vert_F=0$ and $l(x)=0$.

Theorem 4.5 now implies the existence of an $f \in E^{\ast}$ with $f\vert_{(M^{\perp})_{\perp}}=l$. $f$ is in $M^{\perp}$ because of $f\vert_F = l\vert_F =0$ and $M\subseteq F$. But $f(x)\neq 0$ - Contradiction! 
\end{proof}

\setcounter{theo}{9}  %% Remove this when continuing the script!!!!!!!!!!!
\begin{theo}\label{4:t:isoiso}
Let $E$ be a normed space over $\K$, and $F\sse E$ a linear subspace.
\begin{enumerate}
\item[(i)] The linear operator
$$
\Phi : E^*/F^\perp\to F^*,\quad \Phi(f + F^\perp) = f|_F,\;f\in E^*,
$$
is an isometric isomorphism.
\item[(ii)] If $F$ is closed, then the linear operator
$$
\Phi : (E/F)^*\to F^\perp,\quad (\Phi f)(x) = f(x+F),\;x\in E,\,f\in (E/F)^*,
$$
is an isometric isomorphism.
\end{enumerate}
\end{theo}
\begin{proof}
(i). Consider the map $T : E^*\to F^*$, $Tf := f|_F$, $f\in E^*$. We have $\ker T = F^\perp$. Hence, $\Phi$ is well-defined, linear, and injective. By Theorem 4.5\marginpar{Reference!}, for each $\ell\in F^*$ there exists some $f\in E^*$ with $f|_F = \ell$. Hence, $\Phi$ is surjective. Finally, let $f\in E^*$, and choose $g\in E^*$ such that
$$
\|g\| = \|f|_F\|\quad\text{and}\quad g|_F = f|_F.
$$
Then, we obtain
$$
\|f + F^\perp\| \le \|f + (g-f)\| = \|g\| = \|f|_F\| = \|\Phi(f + F^\perp)\|.
$$
On the other hand, for all $g\in F^\perp$,
$$
\|\Phi(f + F^\perp)\| = \|f|_F\| = \|(f+g)|_F\|\le \|f + g\|.
$$
Hence, $\|\Phi(f + F^\perp)\|\le \|f + F^\perp\|$.

(ii). First, $\Phi f : E\to\K$, $x\mapsto f(x + F)$, is linear. Since
$$
|(\Phi f)(x)| = |f(x+F)|\le \|f\|\cdot\|x+F\|\le\|f\|\cdot\|x\|,
$$
we have $\Phi f\in E^*$. If $x\in F$, then $(\Phi f)(x) = f(F) = 0$, hence $\Phi f\in F^\perp$. Therefore, indeed $\Phi : (E/F)^*\to F^\perp$. It is obvious that $\Phi$ is linear and injective. To prove surjectivity, let $g\in F^\perp$, and let $f : E/F\to\K$ be defined by
$$
f(x+F) = g(x),\quad x\in E.
$$
The functional $f$ is is well-defined since $F\sse\ker g$. Moreover, $f$ is linear, and for all $x\in E$, $y\in F$ we have
$$
|f(x+F)| = |g(x)| = |g(x+y)|\le\|g\|\cdot\|x+y\|,
$$
which implies $f\in (E/F)^*$. In addition, $\Phi f = g$, and surjectivity is proved.

It remains to show that $\Phi$ is isometric. For this, note that $|(\Phi f)(x)|\le \|f\|\cdot\|x\|$, $x\in E$, implies $\|\Phi f\|\le\|f\|$ for all $f\in (E/F)^*$. On the other hand, for each $\veps > 0$ there exists $x\in E$ with
$$
\|x + F\| = 1\quad\text{and}\quad |f(x + F)|\ge\|f\| - \veps.
$$
Since $1 = \|x + F\| = \inf_{y\in F}\|x + y\|$, there exists $y\in F$ with $\|x + y\|\le 1+\veps$. This implies $\|\frac{x+y}{1+\veps}\|\le 1$ and hence
$$
\left|(\Phi f)\left(\frac{x+y}{1+\veps}\right)\right| = \frac{|f(x+F)|}{1+\veps}\ge\frac{\|f\| - \veps}{1+\veps}.
$$
Thus $\|\Phi f\|\ge \frac{\|f\| - \veps}{1+\veps}$, which yields $\|\Phi f\|\ge\|f\|$. This proves that $\Phi$ is isometric.
\end{proof}

\begin{lemma}
Let $E$ be a normed space over $\K$. For $x\in E$ let $\hat x : E^*\to\K$ be defined by $\hat x(\ell) = \ell(x)$, $\ell\in E^*$. Then $\Lambda_E : E\to (E^*)^*$, $\Lambda_Ex = \hat x$, is an isometric linear operator.
\end{lemma}
\begin{proof}
$\Lambda_E$ is linear, and for $x\in E$ we have
$$
\sup\{|\hat x(\ell)| : \ell\in E^*,\,\|\ell\|=1\} = \sup\{|\ell(x)| : \ell\in E^*,\,\|\ell\|=1\} = \|x\|,
$$
where the last equality follows from Corollary 4.6\marginpar{Reference!}. This proves that indeed $\hat x\in (E^*)^*$ and that $\|\Lambda_E x\| = \|\hat x\| = \|x\|$ for all $x\in E$.
\end{proof}

\begin{definition}
Let $E$ be a normed space over $\K$, and let $\Lambda_E$ be defined as above. Then $\Lambda_E$ is called {\em canonical map} or {\em embedding} of $E$ in $E^{**} := (E^*)^*$. $E$ is {\em reflexive}, if $\Lambda_E$ is surjective. $E^{**}$ is the {\em bi-dual} of $E$, and $\ol{\Lambda_E(E)}$ is the {\em completion} of $E$.
\end{definition}

\begin{rem}
Only Banach spaces can be reflexive. The class of reflexive spaces is a highly important class of Banach spaces. Intriguingly, there exist non-reflexive Banach spaces, which are isometrically isomorphic to their bi-dual. Moreover, notice that finite-dimensional spaces are always reflexive because of $\dim E^{**} = \dim E$.
\end{rem}

\begin{theo}\label{4:t:refl_1}
Let $E$ be a normed space over $\K$.
\begin{enumerate}
\item[(i)] If $E$ is reflexive and $F\sse E$ a closed linear subspace, then $F$ is also reflexive.
\item[(ii)] If $E$ is a Banach space, then $E$ is reflexive if and only if $E^*$ is reflexive.
\end{enumerate}
\end{theo}
\begin{proof}
(i). We have to show that for each $\vphi\in F^{**}$ there exists some $y\in F$ with $\vphi(f) = f(y)$ for all $f\in F^*$. For this, let $\vphi\in F^{**}$ and let $\psi : E^*\to\K$ be defined by $\psi(\ell) := \vphi(\ell|_F)$, $\ell\in E^*$. Since
$$
|\psi(\ell)|\le\|\vphi\|\cdot\|\ell|_F\|\le\|\vphi\|\cdot\|\ell\|,
$$
we have $\psi\in E^{**}$. $E$ being reflexive then implies that there exists $y\in E$ with $\psi(\ell) = \ell(y)$ for all $\ell\in E^*$. Next, towards a contradiction, assume that $y\notin F$. Then there exists some $\ell\in E^*$ with $\ell(y)\neq 0$ and $\ell|_F = 0$. Hence, $0\neq\ell(y) = \psi(\ell) = \vphi(\ell|_F) = 0$. A contradiction. Finally, for $f\in F^*$ there exists some $\ell\in E^*$ with $\ell|_F = f$, hence
$$
\vphi(f) = \vphi(\ell|_F) = \psi(\ell) = \ell(y) = f(y).
$$
This shows that $F$ is reflexive.

(ii). Let $E$ be reflexive. We need to show that for each $u\in E^{***}$ there exists some $f\in E^*$ with $u(\vphi) = \vphi(f)$ for all $\vphi\in E^{**}$. For this, let $u\in E^{***}$, and set $f(x) := u(\hat x)$, $x\in E$. Then $f\in E^*$. Next, let $\vphi\in E^{**}$. Since there hence exists some $x\in E$ with $\hat x = \vphi$, we obtain
$$
u(\vphi) = u(\hat x) = f(x) = \hat x(f) = \vphi(f).
$$
Therefore, $E^*$ is reflexive.

For the converse, let $E^*$ be reflexive. Then, by the above, $E^{**}$ is reflexive. By (i), also $\Lambda_E(E) = \ol{\Lambda_E(E)}$ is reflexive. The claim now follows from the fact that $E$ and $\Lambda_E(E)$ are (isometrically) isomorphic (see Exercise Sheet 6, Exercise 1(ii)).
\end{proof}

\begin{theo}
Let $E$ be a Banach space and $F\sse E$ a closed linear subspace. Then the following are equivalent:
\begin{enumerate}
\item[(i)] $E$ is reflexive.
\item[(ii)] $F$ and $E/F$ are reflexive.
\end{enumerate}
\end{theo}
\begin{proof}
(i)$\Ra$(ii). By (i) and Theorem \ref{4:t:refl_1}(i), also $F$ is reflexive. By Theorem \ref{4:t:refl_1}(ii), $E^*$ is reflexive, hence $F^\perp$ is reflexive. By Theorem \ref{4:t:isoiso}(ii), $F^\perp$ is isometrically isomorphic to $(E/F)^*$ which is therefore also reflexive. By Theorem \ref{4:t:refl_1}(ii), this finally implies that $E/F$ is reflexive.

(ii)$\Ra$(i). Let $\vphi\in E^{**}$. We will again use the isometric isomorphism
$$
\Phi : (E/F)^*\to F^\perp\sse E^*,\quad (\Phi u)(x) = u(x+F),\;u\in (E/F)^*,\,x\in E,
$$
from Theorem \ref{4:t:isoiso}. Then we can define $\psi\in (E/F)^{**}$ by
$$
\psi(u) := \vphi(\Phi u),\quad u\in (E/F)^*.
$$
Since $E/F$ is reflexive, there exists some $x\in E$ with $\wh{x+F} = \psi$, i.e.
$$
\vphi(\Phi u) = \psi(u) = (\wh{x+F})(u) = u(x+F) = (\Phi u)(x) = \hat x(\Phi u),\quad u\in(E/F)^*.
$$
Hence, $(\vphi - \hat x)|_{F^\perp} = 0$.

To utilize the reflexivity of $F$, we next define a suitable $\rho\in F^{**}$. For each $f\in F^*$, choose some $g\in E^*$ with $g|_F = f$ and $\|g\| = \|f\|$. Then define
$$
\rho(f) := (\vphi - \hat x)(g).
$$
This is a proper definition since for two extensions $g,h\in E^*$ of $f$ we have $(g - h)|_F = 0$ and thus $g - h\in F^\perp$. A similar argument shows that $\rho$ is linear. Moreover,
$$
|\rho(f)|\le\|\vphi - \hat x\|\cdot\|g\| = \|\vphi - \hat x\|\cdot\|f\|.
$$
Thus, $\rho\in F^{**}$. As $F$ is reflexive, there exists some $y\in F$ with $\rho(f) = f(y)$ for all $f\in F^*$. Now, we conclude that for all $g\in E^*$ we have
$$
\hat y(g) = g(y) = (g|_F)(y) = \rho(g|_F) = (\vphi - \hat x)(h)
$$
with some $h\in E^*$ satisfying $h|_F = g|_F$ and $\|h\| = \|g|_F\|$. Hence, $h - g\in F^\perp$ and thus $\hat y(g) = (\vphi - \hat x)(g)$ for all $g\in E^*$. Equivalently, $\vphi = \hat x + \hat y = \wh{x + y}\in\Lambda_E(E)$, which shows that $E$ is reflexive.
\end{proof}

\begin{theo}
Let $E$ and $F$ be normed spaces, $E\neq\{0\}$. If $\calL(E,F)$ is complete, then so is $F$.
\end{theo}
\begin{proof}
First, choose $x_0\in E$ with $\|x_0\|=1$. Then there exists some $f\in E^*$ with $f(x_0) = \|x_0\| = 1 = \|f\|$. Next, let $(y_n)_{n\in\N}\sse F$ be a Cauchy sequence, and define $T_n : E\to F$ by
$$
T_nx := f(x)y_n,\quad x\in E.
$$
Since $\|T_nx\|\le\|f\|\|x\|\|y_n\| = \|y_n\|\|x\|$ for $x\in E$, we have $T_n\in\calL(E,F)$. Further, $\|T_n x- T_m x\| = |f(x)|\|y_n - y_m\|\le\|y_n - y_m\|\|x\|$ implies
$$
\|T_n - T_m\|\le\|y_n - y_m\|.
$$
Hence $(T_n)_{n\in\N}$ is a Cauchy sequence in $\calL(E,F)$ and thus converges to some $T\in\calL(E,F)$. This implies $y_n = T_nx_0\to Tx_0$ as $n\to\infty$, i.e.\ $(y_n)_{n\in\N}$ converges in $F$.
\end{proof}


\section{The open mapping, closed graph and Banach-Steinkraus theorem}

\begin{lemma}
	Let $E$ be a normedm $F$ a Banach space and $T \in \calL(E,F)$ surjective. Then
	\begin{equation*}
		K_r(0_F) \subset \closure{T(K_r(0_E))} \qquad \text{ for some } r > 0
	\end{equation*}
\end{lemma}

%thrm 7
\begin{theo}[Uniform Boundedness Priniple]
	\label{5:t:ubp}
	Let $E$ be a Banach space, $F$ a normed space and let $\calT\in\calL(E,F)$. Let $\calT$ be pointwise bounded, i.e., for each $x\in E$ there exists $M_x<\infty$ sucht that $\|Tx\|\leq M_x\quad\forall T\in\calT$. Then $\calT$ is bounded, i.e., there exists some $M<\infty$ sucht that $\|T\|<M\quad\forall T\in\calT$.
\end{theo}
\begin{proof}For $\nin$ let \[E_n=\{x\in E:\|Tx\|\leq m\quad\forall T\in\calT\}.\]
By hypothesis, $E=\bigcup\limits_{n=1}^\infty E_n$.
Let $x=\lim\limits_{j\to\infty}x_j$ with $x_j\in E_n$ for fixed $n$.
Since $\|Tx_j\|\leq n$ for all $j$ we have \[\|Tx\|=\lim_{j\to\infty}\|Tx_j\|\leq m.\]
Thus $x\in E_n$, and hence $E_n$ is closed.
By remark \ref{rem:baire} on Baire's Theorem, there exists some $n_0\in\N$ with \[\mathring E_{n_0}\neq\emptyset.\]
Hence, $K_r(x)\subseteq E_{n_0}$ for some $x\in E_n$, $r>0$.
Let $y\in E$ with $\|y\|\leq r$.
Then $y+x\in K_r(x)=x+K_r(0)$.
This implies \begin{align}\begin{split}\label{eqT5.8}\|Ty\|&=\|T(y+x)-Tx\|\\&\leq\|T(y+x)\|+\|Tx\|\leq2n_0\quad\forall T\in\calT.\end{split}\end{align} Now let $y\in E$, $y\neq0$ arbitrary. Then \[\begin{array}{rrcl}&\frac r{\|y\|}\|Ty\|&=&\left\|T\left(\frac{ry}{\|y\|}\right)\right\|\overset{\ref{eqT5.8}}\leq 2n_0\\\Ra&\|Ty\|&\leq&\frac{2n_0}r\|y\|\\\Ra&\|T\|&\leq&\frac{2n_0}r.\end{array}\]
\end{proof}
\begin{bsp}In general, Theorem \ref{5:t:ubp} does not hold, if $E$ is \emph{not} a Banach space. Example: \begin{align*}E&=\{x=(x_n)_\nin\in\ell_1:x_n=0\text{ for almost all }\nin\}\\F&=\K\\f_n(x)&=nx_n\quad\forall x\in E,\nin\\\calT&=\{f_n:\nin\}.\end{align*}
We see that $\calT$ is pointwise bounded, since $x_n=0$ from some $n\geq N$ on.

BUT $\|f_n\|=n,\,\nin$.
\end{bsp}
\begin{kor}\label{5:c:10}
Let $E$ be a Banach Space, $F$ a normed space and $T_n\in\calL(E,F)$.
Suppose for every $x\in E$, $(T_nx)_\nin$ is convergent in $E$.
Then define $T\colon E\to F$ by \[Tx:=\lim_{n\to\infty}T_nx.\]
Then $T\in\calL(E,F)$, $(\|T_n\|)_\nin$ is bounded, and \[\|T\|\leq\liminf_{n\to\infty}\|T_n\|.\]
\end{kor}
\begin{proof} By definition, $T$ obviously is linear. $(\|T_nx\|)_\nin$ is bounded. By Theorem \ref{5:t:ubp}, $\|T_n\|\leq M$ for all $\nin$.
Hence, for all $x\in E$ \[\|Tx\|=\lim_{n\to\infty}\|T_nx\|\leq M\|x\|.\]
This shows that $T\in\calL(E,F)$.
Let $(\|T_{n_k}\|)_{k\in\N}$ be a convergent subsequence of $(\|T_n\|)_\nin$.
Then \[\|Tx\|=\lim_{k\to\infty}\|T_{n_k}\|\leq\|x\|\lim_{k\to\infty}\|T_{n_k}.\]
Thus \[\|T\|\leq\lim_{k\to\infty}\|T_{n_k}\|,\] and hence \[\|T\|\leq\liminf_{n\to\infty}\|T_n\|.\]
\end{proof}
\begin{lemma}\label{5:l:11}
Let $E$ be a normed space and $F$ be a Banach space.
Then let $E_0$ be a dense linear subspace of $E$, and $T_0\in\calL(E_0,F)$.
There exists a unique $T\in\calL(E,F)$ with \[T\vert_{E_0}=T_0\quad\text{and}\quad\|T\|=\|T_0\|.\]
\end{lemma}
\begin{proof}Follows with Lemma 3.6\marginpar{Reference!}.
\end{proof}
\begin{theo}[Banach-Steinhaus Theorem]\label{5:t:bst}\hfill

\begin{enumerate}[(i)]\item Let $E$ be a Banach space and $F$ a normed space. Further, let $T_n\in\calL(E,F)$, $\nin$.
If $(T_n)_\nin$ is pointwise convergent to some $T\colon E\to F$ which is linear, then \[\sup_\nin\|T_n\|<\infty.\]

\item Let $E$ be a normed space and $F$ a Banach space. Further, let $T\in\calL(E,F)$, $\nin$. If
\begin{enumerate}[(a)]\item $\sup\limits_\nin\|T_n\|<\infty$
\item there exists a dense linear subspace $E_0$ of $E$ such that $(T_nx)_\nin$ is convergent in $F$ for each $x\in E_0$,
\end{enumerate}
then there existst some $T\in\calL(E,F)$ with \[Tx=\lim_{n\to\infty}T_nx\quad\text{for all }x\in E.\]
\end{enumerate}
\end{theo}
\begin{proof}\hfill

\begin{enumerate}[(i)]\item This is Cor. \ref{5:c:10}
\item For each $y\in E_0$ set \[T_0y:=\lim_{n\to\infty}T_ny.\quad\text{(exists by (b))}\]
$T_0$ is linear and \[\|T_0y\|=\lim_{n\to\infty}\|T_ny\|\leq\ub{\sup_\nin\|T_n\|\|y\|}_{<\infty\text{ by (a)}}.\]
Hence $T_0\in\calL(E,F)$. By Lemma \ref{5:l:11}, there exists some $T\in\calL(E,F)$ with $T\vert_{E_0}=T_0$.
Let $x\in E$, and $\veps>0$. Then let $y\in E_0$ with \[\|x-y\|\leq\veps\text{ and }N\in\N\text{ with }\|T_ny-T_0y\|\leq\veps\ \forall n\geq N.\quad\text{(b)}\]
Then for all $n\geq N$, \begin{align*}\|T_nx-Tx\|&\leq\ub{\|T_nx-T_ny\|}_{\leq\|T_n\|\cdot\|x-y\|}+\ub{\|T_ny-T_0y\|}_{\leq\veps}+\ub{\|T_0y-Tx\|}_{\leq\|T\|\cdot\|y-x\|}\\&\leq\|T_n\|\|x-y\|+\veps+\|T\|\|y-x\|\\&\leq\veps(\|T_n\|+1+\|T\|)\\&\leq\veps\biggl(\ub{\sup_\nin\|T_n\|}_{<\infty\text{ by (a)}}+1+\|T\|\biggr).\end{align*}
This implies $T_nx\to Tx$ as $n\to\infty$.
\end{enumerate}
\end{proof}
\begin{theo}\label{5:t:13}
Let $E$ be a normed space and $M\subseteq E$.
Then the following conditions are equivalent:
\begin{enumerate}[(i)]\item $M$ ist bounded.
\item $f(M)\subseteq\K$ is bounded for all $f\in E^*$.
\end{enumerate}
\end{theo}
\begin{bem}[Geometric interpretation of \ref{5:t:13}]
Suppose that for every closed hyperplane $H$ in $E$ (kernel of $f$) some $c$ exists with $M$ between $H+c$ and $H-c$.
Then $M$ is already contained in a ball. \hfill ((ii)$\Ra$(i))
\end{bem}
\begin{proof}\hfill

\begin{itemize}\item[(i)$\Ra$(ii)] This follows from \[\|x\|<c\Ra\|f(x)\|\leq c\|f\|\quad\forall x\in M.\]
\item[(ii)$\Ra$(i)] Consider the set \[\hat M=\{\hat x\colon f\mapsto f(x):x\in M\}\subseteq\calL(E^*,\K)=E^{**}.\]
Since $\hat M(f)=f(M)$, by (ii) $M$ is pointwise bounded. By Theorem \ref{5:t:ubp}, $\hat M$ is bounded. Since $\hat{}$ is isometric, also $M$ is bounded.
\end{itemize}
\end{proof}
\begin{kor} Let $E,F$ be normed spaces and $T\colon E\to F$ be linear. Then \[T\in\calL(E,F)\eq f\circ T\in E^*\text{ for all }f\in F^*.\]
\end{kor}
\begin{proof}\hfill

\begin{align*}
	{rcl}T\text{ is bounded}&\eq&T(K_1(0))\text{ is bounded in }F\\&\overset{\ref{5:t:13}}\eq&f(T(K_1(0)))\text{ is bounded for all }f\in F^*\\&\eq&f\circ T\in E^*\text{ for all }f\in F^*.
\end{align*}
\end{proof}


\newpage
\begin{thebibliography}{}
\bibitem{Book1} H. W. Alt, \emph{Funktionalanalysis. Eine anwendungsorientierte Einf�hrung}, 6. Aufl., Springer, 2011.
\bibitem{Book2} J. B. Conway, \emph{A course in functional analysis}, 2. Aufl., Springer, 1990.
\bibitem{Book3} H. Heuser, \emph{Funktionalanalysis. Theorie und Anwendung}, 4. Aufl., Teubner, 2006
\bibitem{Book4} W. Rudin, \emph{Functional Analysis}, McGraw-Hill, 1991.
\bibitem{Book5} D. Werner, \emph{Funktionalanalysis}, 7. Aufl., Springer, 2011.
\bibitem{Book6} O. Christensen, \emph{Frames and bases: an introduction course}, Birkh�user, 2008

\end{thebibliography}
\addcontentsline{toc}{section}{Literatur}
\end{document}
